% concatenated from QPL2026.tex
\documentclass[submission,copyright,creativecommons]{eptcs}
\providecommand{\event}{SOS 2007} % Name of the event you are submitting to

\usepackage{iftex}
\usepackage{todonotes}

\newcommand{\marn}[1]{\todo[color=green!40]{#1}}
\newcommand{\note}[1]{\todo[inline, color=gray!20]{#1}}
\newcommand{\warn}[1]{{\color{red}#1}}

\ifpdf
  \usepackage{underscore}         % Only needed if you use pdflatex.
  \usepackage[T1]{fontenc}        % Recommended with pdflatex
\else
  \usepackage{breakurl}           % Not needed if you use pdflatex only.
\fi

\usepackage{enumerate,xspace}
\usepackage{amsmath,amssymb,wasysym}
\usepackage[all]{xy}
\usepackage{multicol}
\usepackage{proof}
\usepackage{tikz-cd}

\newtheorem{observation}{Remark}[section]
\newtheorem{lemma}[observation]{Lemma}  %%share counter with remark
\newtheorem{theorem}[observation]{Theorem}
\newtheorem{definition}[observation]{Definition}
\newtheorem{example}[observation]{Example}
\newtheorem{remark}[observation]{Remark}
\newtheorem{notation}[observation]{Notation}
\newtheorem{proposition}[observation]{Proposition} 
\newtheorem{corollary}[observation]{Corollary} 
\newtheorem{conjecture}[observation]{Conjecture}
\newcommand{\proof}{\noindent{\sc Proof:}\xspace}
\def\endproof{~\hfill$\Box$\vskip 10pt}

\title{Generalized Inverses of Quantum Channels: \\a categorical perspective}
\author{Robin Cockett\thanks{Partially funded by NSERC.}
\institute{Department of Computer Science\\
University of Calgary, Canada}
\email{robin@ucalgary.ca}
\and
Jean-Simon Pacaud Lemay%\thanks{Financially supported by...}
\institute{School of Mathematical and Physical Sciences\\
Macquarie University, Australia}
\email{js.lemay@mq.edu.au}
\and
Priyaa Varshinee Srinivasan
\thanks{This work was co-funded by the European
Union and Estonian Research Council through the Mobilitas 3.0 (MOB3JD1227).}
\institute{Department of Software Sciences\\
Tallinn University of Technology, Estonia}
\email{priyaavarshinee@gmail.com}
}

\def\event{QPL 2026}
\def\titlerunning{Generalized Inverses of Quantum Channels: a categorical perspective}
\def\authorrunning{J.R.B. Cockett, J.-S. P. Lemay, \& P. V. Srinivasan}

\begin{document}
\maketitle
% Robin doing his bit here!!!!  Actually I have done corrections here!
%Great thanks. Did you make those corrections 2-3 days ago? if yes I can see those in the history

\begin{abstract}
A quantum channel is defined as being completely positive (CP) and trace preserving (TP). While not every quantum channel is invertible or reversible, every quantum channel admits various kinds of generalized inverses such as the Moore-Penrose inverse and the Drazin inverse. A generalized inverse of a quantum channel may not itself be a quantum channel: it often fails to be CP.  However, generalized inverses still play an important role in quantum error mitigation.  Here, because it is often desirable for the generalized inverse of a quantum channel to be at least TP, the Drazin inverse, which is TP, is favoured over the Moore-Penrose inverse, which is not in general TP. 

In this paper, we take a categorical perspective on generalized inverses of quantum channels. This allows us to give a simple proof of the fact that the Drazin inverse of a quantum channel is always TP. It also allows us to show that for unital quantum channels, the Drazin inverse is also unital. We then generalize this result to dagger Drazin inverses, which allows us to show that for unital quantum channels, the Moore-Penrose inverse is both TP and unital as well. This opens the door to new applications of both the Drazin inverse and Moore-Penrose inverse in quantum information theory and, in particular, in quantum error mitigation. 

%We also consider various examples of a quantum channel in a dagger traced symmetric monoidal closed category, including some whose generalized inverse is again a quantum channel. \textcolor{purple}{* I think it is TP and not a quantum channel, because generalized inverse of quantum channel is not necessarily CP.}

\end{abstract}

\section{Introduction}

%Quantum channels are one of the fundamental notions in quantum information. Mathematically, if $H$ is a finite dimensional Hilbert space and $\mathcal{L}(H)$ is the space of complex linear operators on $H$, then a quantum channel \cite[Chap 2.2]{watrous2018theory} is a complex linear map $\Phi: \mathcal{L}(H) \to \mathcal{L}(H^\prime)$ which is (i) completely positive (CP), meaning for any other finite dimensional Hilbert space $K$, $\Phi \otimes \mathsf{id}_{\mathcal{L}(K)}$ maps positive semidefinite operators to positive definite operators, and (ii) trace preserving (TP), meaning that $\mathsf{Tr}(\Phi(f)) = \mathsf{Tr}(f)$ for all $f \in \mathcal{L}(H)$. These conditions insure that quantum channels preserve quantum states, i.e. density operators. As complex linear maps, we can ask what algebraic properties quantum channels satisfy, which in turn help gives us better understanding of our quantum channels. 

%In particular, one can ask whether or not a quantum channel has an inverse. Of course, as with linear maps, not all quantum channels have an inverse, and even if one does, the inverse may not be a quantum channel, often failing to be CP. On the other hand, all complex linear maps have various kinds of \textit{generalized} inverses \cite{campbell2009generalized}. Thus one can consider various kinds of generalized inverses of quantum channels, such as the famous Moore-Penrose inverse \cite[Chap 1]{campbell2009generalized} or even the Drazin inverse \cite[Chap 7]{campbell2009generalized}. That said, while every quantum channel admits generalized inverses, these generalize inverses need not be quantum channels, often failing to be CP or TP, or even both. Nevertheless, generalized inverses of quantum channels still have useful applications. 

In the era of noisy intermediate-scale quantum (NISQ) computing \cite{Pre18}, both quantum error correction \cite{chatterjee2023quantum} and quantum error mitigation \cite{cao2021nisq,cai2023quantum} are essential for reliable quantum computation. Generalized inverses \cite{campbell2009generalized} of quantum channels (completely positive (CP) trace-preserving (TP) linear maps) play an important role in quantum error mitigation. 

While quantum error correction protocols  aim to eliminate noise by applying an inverse quantum channel, quantum error mitigation protocols  reduce noise by post-processing measurement results using generalized inverses  when the noise  is not reversible. For quantum error mitigation protocols, the generalized inverses of quantum channels no longer need to be physically implementable, hence need not be CP. However, it is still desirable for the generalized inverse to be TP since, for example, it can provide better stability for computer simulations. 
%Otherwise, probabilistic error cancellation would require explicit renormalization which may introduce systematic bias in the expected values.
For this reason, the Drazin inverse is favoured over the Moore-Penrose inverse for quantum error mitigation since the Drazin inverse \cite[Chap 7]{campbell2009generalized} of a TP map is again TP \cite[Thm 2]{cao2021nisq}, while the Moore-Penrose inverse \cite[Chap 1]{campbell2009generalized} need not be \cite[Ex 2]{cao2021nisq}. 

%------

%Quantum error correction protocols \cite{chatterjee2023quantum} aim to eliminate the effects of noise by applying a reverse quantum channel. When such a reverse quantum channel is not available, quantum error mitigation\cite{cao2021nisq,cai2023quantum} protocols attempt to reduce noise in quantum computation by post-processing measurement results using a \textit{generalized} inverse (a.k.a. pseudoinverse) such as Moore-Penrose\cite[Chap 1]{campbell2009generalized} or Drazin inverse\cite[Chap 7]{campbell2009generalized}. Unlike an inverse of a quantum channel, a generalized inverse may not itself be a quantum channel: failing to either be completely positive (CP), trace preserving (TP), or even both, hence not physically implementable.

%While error mitigation protocols does not require the generalized inverse of a quantum channel to be CP, it is still desirable that the inverse is TP, in particular since it can provide better stability for computer simulations. Otherwise, probabilistic error cancellation would require explicit renormalization which may introduce systematic bias in the expected values. As such for quantum error mitigation, the Drazin inverse is favoured over the Moore-Penrose inverse since the Drazin inverse of a TP map is again TP \cite[Thm 2]{cao2021nisq}, while the Moore-Penrose inverse need not be \cite[Ex 2]{cao2021nisq}.

%---------------------

%Complex linear maps between finite dimensional Hilbert spaces have various kinds of \textit{generalized} inverses, such as the famous Moore-Penrose inverse \cite[Chap 1]{campbell2009generalized} or even the Drazin inverse \cite[Chap 7]{campbell2009generalized}. As such, one can consider these generalized inverses of quantum channels. However, a generalized inverse of quantum channel may not itself be a quantum channel: failing to either be completely positive (CP), trace preserving (TP), or even both. Nevertheless, generalized inverses of quantum channels play an important role in quantum error mitigation \cite{cao2021nisq,cai2023quantum}.

%While quantum error correction protocols \cite{chatterjee2023quantum} aim to eliminate the effects of noise by applying a reverse quantum channel, quantum error mitigation protocols attempt to reduce noise in quantum computation by post-processing measurement results using generalized inverses, when a reverse quantum channel is not available. As such, for quantum channels for quantum error mitigation protocols, the generalized inverses no longer need to be CP. However, it is still desirable for the generalized inverse to be TP, in particular since it can provide better stability for computer simulations. Otherwise, probabilistic error cancellation would require explicit renormalization which may introduce systematic bias in the expected values. As such for quantum error mitigation, the Drazin inverse is favoured over the Moore-Penrose inverse since the Drazin inverse of a TP map is again TP \cite[Thm 2]{cao2021nisq}, while the Moore-Penrose inverse need not be \cite[Ex 2]{cao2021nisq}. 

Quoting \cite{cao2021nisq}: ``{\em To our knowledge, generalized non-CPTP inverses for non-invertible quantum channels have not been studied extensively in previous literature. The new understanding of the natural representation opens the possibility of studying the structures and properties of these maps from mathematical interests. Different constructions of generalized inverse maps bring new properties. This research direction can provide a guideline for implementing these maps in EM and also in other areas of quantum information sciences.}" Motivated by this philosophy, the purpose of this paper is to continue the study generalized inverses of quantum channels from the category theory perspective, combining categorical quantum foundations and the study of generalized inverses in categories. Indeed, abstractions of quantum channels, especially the notion of being CP, are well-established and well-studied concepts in dagger compact closed categories \cite{coecke2016_cpinf,Coecke_2012,coecke2016categories,Cunningham_2015,Gogioso_2019,selinger2007dagger}. On the other hand, there has been recent interest on Moore-Penrose and Drazin inverses in dagger categories \cite{Cockett2023Moore,cockett2025dagger,cockett2024drazin,EPTCS426.3}. 

The main theoretical result of \cite{cao2021nisq} is that the Drazin inverse of a TP map is again TP, the proof \cite[App A]{cao2021nisq} of which involved  heavy duty linear algebra. The first main consequence of our categorical approach is a much more streamlined and shorter proof of the same result using two simple commuting triangles  (Prop \ref{prop:DrazinTP}), which are a consequence of the algebraic properties of Drazin inverses available in any category \cite[Prop 3.10]{cockett2024drazin}. This also allows us to show that for \textit{unital} quantum channels, the Drazin inverse is also unital (Lemma \ref{lemma:DrazinU}). 

One disadvantage of Drazin inverses is that they can only be defined for endomaps. This is solved when one considers dagger Drazin inverses, a generalization of Drazin inverses in dagger categories \cite{cockett2025dagger}. We provide the analogue of \cite[Prop 3.10]{cockett2024drazin} for dagger Drazin inverses (Prop \ref{dagger-Drazin-commuting}), which then allows us to show that the dagger Drazin inverse of a TP and untial map is again TP and unital (Prop \ref{dag-Drazin-TP}). 

As mentioned above, one of the key motivating observations in \cite{cao2021nisq} was that while the Drazin inverse of a quantum channel is always TP, its Moore-Penrose inverse need not be. The second main result of our categorical approach is that we can show that for unital quantum channels, the Moore-Penrose inverse is both TP and unital. In fact, we show that a map is TP and unital if and only if its Moore-Penrose inverse is TP and unital (Prop \ref{MP-TP}). This observation can potentially lead to new applications of Moore-Penrose inverses in quantum information theory, and in particular to applications of the Moore-Penrose inverse of unital quantum channels in quantum error mitigation. For example, mixed unitary channels are central to the study of noise and decoherence in quantum systems with wide-ranged applications including in quantum cryptography \cite{ambainis2000private,ambainis2004small,hayden2004randomizing} and error mitigation \cite{TuJ25}. Since a mixed unitary channel is a unital quantum channel, as a direct application of our result, we know that its Moore-Penrose inverse is both TP and unital. To our knowledge, this behaviour of mixed unitary channels has not been established previously. 

%While quantum channels are often generalized in the setting of dagger compact closed categories, for the story of this paper, we work in the slightly more general framework of dagger traced symmetric monoidal closed categories. While dagger compact closed categories are very well behaved and admit a very practical graphical calculus \cite{selinger2007dagger}, we find the setting of dagger traced symmetric monoidal closed categories ($\dagger$-TMSCC) slightly more intuitive for the story of generalized inverses, in particular since it places emphasis on the trace, which is the main structure we wish to preserve with our generalized inverses. As it is always useful to provide new points of view, we hope that this dagger traced symmetric monoidal closed category  perspectives will provide new insights on quantum channels -- though we leave the rigorous study of CP maps in this setting for future considerations. All our examples of dagger traced monoidal closed categories are dagger compact closed. We leave it as a future work to investigate this setting in greater detail. 

\textbf{Acknowledgements:} The authors thank Chris Heunen, Masahito Hasegawa, Cole Comfort, Amar Hadzihasanovic, and Martti Karvonen for useful discussions and answering our multiple questions. 

\section{Quantum Channels in Dagger Compact Closed Categories}

In this section, we review the notion of abstract quantum channels of dagger compact closed categories. We assume that the reader is familiar with the basics of dagger monoidal categories. For a detailed introduction to dagger compact closed categories, we refer the reader to \cite{selinger2010survey} (which also includes the graphical calculus for dagger compact closed categories, which we will not need to make use of in this paper), and for a detailed introduction to complete positivity in this setting, we refer the reader to \cite{selinger2007dagger}. 

For an arbitrary category $\mathbb{X}$, we will denote objects using capital letters $A,B,C$, etc. and maps will denoted as arrows $f: A \to B$, with homsets as $\mathbb{X}(A,B)$. Identity maps will denoted as $1_A: A \to A$ (or simply $1$ when there is no confusion) while composition will be denoted diagrammatically, so the composite of $f: A \to B$ and $g: B \to C$ is written as $f;g: A \to C$, which is, first do $f$ then $g$. For simplicity, we allow ourselves to work in a \textit{strict} monoidal category setting, meaning that the associativity and unital isomorphisms of the monoidal product are identities. So, for a symmetric (strict) monoidal category $\mathbb{X}$, we denotes its monoidal product as $\otimes$, its monoidal unit as $I$, and the natural symmetry isomorphism as $\sigma_{A,B}: A \otimes B \to B \otimes A$. Strictness allows us to write down $A_1 \otimes \hdots \otimes A_n$ and $A \otimes I = A = I \otimes A$. 

Now recall that a \textbf{compact closed category} \cite[Sec 4.8]{selinger2010survey} is a symmetric monoidal category such that every object has a chosen dual. Explicitly, for every object $A$, there is an object $A^\ast$, called the \textbf{dual} of $A$, equipped with a pair of maps $\cup_A: A^\ast \otimes A \to I$, called the \textbf{cup}, and $\cap_A: I \to A \otimes A^\ast$, called the \textbf{cap}, such that the famous snake equations hold. For convenience, we will work with the convention that $I^\ast = I$ and $(A \otimes B)^\ast = B^\ast \otimes A^\ast$. On the other hand, a \textbf{dagger category} \cite[Sec 7]{selinger2010survey} is a category $\mathbb{X}$ equipped with a contravariant functor $\dagger: \mathbb{X} \to \mathbb{X}$ which is the identity on objects and also involutive. Explicitly, a dagger associates every map $f: A \to B$ to a map of dual type $f^\dagger: B \to A$, called its \textbf{adjoint}, such that (i) $(f;g)^\dagger= g^\dagger; f^\dagger$, (ii) $1^\dagger = 1$, and (iii) $f^{\dagger \dagger} = f$. A \textbf{dagger symmetric monoidal category} \cite[Sec 7.2]{selinger2010survey} is a symmetric monoidal category which has a dagger which is compatible with the symmetric monoidal structure in the sense that $\otimes$ is a dagger functor, which amounts to asking that $(f \otimes g)^\dagger = f^\dagger \otimes g^\dagger$, and also that the adjoint of the natural symmetry isomorphism is its inverse, so we have $\sigma_{A,B}^\dagger = \sigma^{-1}_{A,B} = \sigma_{B,A}$ (in other words, $\sigma$ is unitary). Then a \textbf{dagger compact closed category} \cite[Sec 7.4]{selinger2010survey} (or $\dagger$KCC for short) is a compact closed category which is also a dagger symmetric monoidal category such that the adjoint of the cap (resp. cup) is the cup (resp. cap) up to a twist, so $\cap_A^\dagger =  \sigma_{A,A^\ast}; \cup_A$ and $\cup_A^\dagger = \cap_A; \sigma_{A, A^\ast}$.

\begin{example} Let $\mathbb{C}$ be the field of complex numbers, and let $\mathsf{FHILB}$ be the category of finite dimensional (complex) Hilbert spaces and $\mathbb{C}$-linear maps between them. Then $\mathsf{FHILB}$ is a $\dagger$KCC. The monoidal structure is given by the usual tensor product of Hilbert spaces, $H \otimes K$, and the monoidal unit is $\mathbb{C}$. The dual of $H$ is the space of linear functionals, $H^\ast = \mathsf{FHILB}(H, \mathbb{C})$. In bra-ket notation, if $\lbrace \vert e_i \rangle \rbrace^n_{i=1}$ is a basis of $H$, with $\lbrace \langle e_i \vert \rbrace^n_{i=1}$ the dual basis for $H^\ast$, the cup and cap are defined as $\cup_H\left(\langle e_i \vert \otimes \vert e_j \rangle \right) = \langle e_i \vert e_j \rangle$ and $\cap_H(1) = \sum^n_{i=1} \vert e_i \rangle\langle e_i \vert$. The dagger is given by the adjoint, that is, for a $\mathbb{C}$-linear map $f: H \to K$, $f^\dagger: K \to H$ is the unique map such that $\langle f(x), y \rangle = \langle x, f^\dagger(y)\rangle$. %Lastly, the biproduct is given by the direct sum of Hilbert spaces, $H \oplus K$, which in turn induces the usual sum of $\mathbb{C}$-linear maps. 
\end{example}

\begin{example} Let $\mathsf{REL}$ be the category of sets and relations, where recall the objects are sets and where a map $R: X \to Y$ is a subset $R \subseteq X \times Y$. Then $\mathsf{REL}$ is a $\dagger$KCC. The monoidal structure is given by the Cartesian product of sets, $X \times Y$, and the monoidal unit is a chosen singleton $\lbrace \ast \rbrace$. The dual of $X$ is itself, $X^\ast$, while the cup and cap are the subsets $\cap_X = \lbrace (\ast, x,x) \vert \forall x \in X \rbrace$ and $\cup_X: \lbrace (x,x,\ast) \vert \forall x \in X\rbrace$. The dagger is given by the converse relation, so for $R \subseteq X \times Y$, $R^\dagger = \lbrace (y,x) \vert (x,y) \in R \rbrace \subseteq Y \times X$. %Lastly, the biproduct is given by the disjoint unit of sets, $X \sqcup Y$, while the induced sum is given by the union, $R \cup S$. 
\end{example}

Now before we turn our attention to reviewing quantum channels in a $\dagger$KCC, for the intuition of the story of this paper, we find it useful to consider $\dagger$KCC from the perspective of being \textit{closed} and having a \textit{trace}. This will help provide some useful notation that may be more familiar for certain readers. 

So briefly recall that a \textbf{symmetric monoidal closed category} is a symmetric monoidal category $\mathbb{X}$ such that for every pair of objects $A$ and $B$, there is an object $[A,B]$ called the \textbf{internal hom}, such that we have natural isomorphisms $\mathbb{X}(C \otimes A, B) \cong \mathbb{X}(C, [A,B])$. Intuitively we think of $[A,B]$ as an object representing all maps from $A$ to $B$. Now every compact closed category is also symmetric monoidal closed by setting $[A,B] = A^\ast \otimes B$, in fact compact closed categories are precisely the symmetric monoidal closed categories such that the canonical map $[A,I] \otimes B \to [A,B]$ is an isomorphism. Of particular importance for the story of quantum channels will be the object $[A,A]= A^\ast \otimes A$, which we think as the object of endomaps on $A$. 

On the other hand, a \textbf{traced symmetric monoidal closed category} \cite{hasegawa2009traced} is a symmetric monoidal closed category equipped with a \textbf{trace operator} $\mathsf{Tr}$, which is a family of functions indexed by triples of objects, $\mathsf{Tr}^X_{A,B}: \mathbb{X}(X \otimes A, X \otimes B) \to \mathbb{X}(A,B)$, where for a map $f: X \otimes A \to X \otimes B$, the resulting map $\mathsf{Tr}^X_{A,B}(f): A \to B$ is called the \textbf{partial trace of $f$ at $X$}. This trace operator $\mathsf{Tr}$ also satisfies well-known axioms \cite[Sec 5]{selinger2010survey} such as naturality in the $A$ and $B$ arguments, extranaturality in the $X$ argument, superposition, and yanking. Now in the special case that $A=B=I$, we will simply denote the trace operator as $\mathsf{Tr}^X = \mathsf{Tr}^X_{I,I}$, and in this case, for an endomap $f: X \to X$, we call $\mathsf{Tr}^X(f): I \to I$ its \textbf{trace}. In a traced symmetric monoidal closed category, we can internalize the trace operator to obtain maps of type $\mathsf{tr}^X_{A,B}: [X \otimes A, X \otimes B] \to [A,B]$ and $\mathsf{tr}^X = \mathsf{tr}^X_{I,I}: [X,X] \to I$ respectfully. 

Now every compact closed category is a traced symmetric monoidal closed category in a unique way \cite[Sec 3.2]{hasegawa2009traced}. In particular, the internal trace operators are given precisely by the cups, so $\mathsf{tr}^X_{A,B} = 1_{A^\ast} \otimes \cup_X \otimes 1_B$ and $\mathsf{tr}^X= \cup_X$. 
%\begin{align}
 %   \mathsf{tr}^X_{A,B} = 1_{A^\ast} \otimes \cup_X \otimes 1_B && \mathsf{tr}^X= \cup_X
%\end{align}
Now if the cups gives the internal trace, what do the caps give? Caps instead give the internal way of picking out identity maps in the internal hom. Indeed, in a symmetric monoidal closed category, we have canonical maps $\mathsf{u}^X : I \to [X,X]$ and $\mathsf{u}^X_{A,B}: [A,B] \to [X \otimes A, X \otimes B]$, where the former is interpreted as picking out $1_X$ in $[X,X]$, while the latter sends $f: A \to B$ to $1_X \otimes f$. In a compact closed category, these maps are given by the cap (up to a twist): $ \mathsf{u}^X_{A,B} =  (1_{A^\ast} \otimes \cap_X \otimes 1_B);(1_{A^\ast} \otimes \sigma_{X, X^\ast} \otimes 1_B)$ and $\mathsf{u}^X= \cap_X; \sigma_{X, X^\ast}$. 
%\begin{align}
  %  \mathsf{u}^X_{A,B} =  (1_{A^\ast} \otimes \cap_X \otimes 1_B);(1_{A^\ast} \otimes \sigma_{X, X^\ast} \otimes 1_B)  && \mathsf{u}^X= \cap_X; \sigma_{X, X^\ast}
%\end{align}
Furthermore, in a $\dagger$KCC, we also have that $({\mathsf{u}^X_{A,B}})^\dagger = \mathsf{tr}^X_{A,B}$ and $({\mathsf{u}^X})^\dagger = \mathsf{tr}^X$. 
%the following equalities: 
%\begin{align}\label{dagger-tr-u}
 %  {\mathsf{u}^X_{A,B}}^\dagger = \mathsf{tr}^X_{A,B} &&  {\mathsf{u}^X}^\dagger = \mathsf{tr}^X
%\end{align}

Lastly, the final ingredient we need to properly define is the ability to conjugate maps. In a symmetric monoidal closed category, for every map $f: A \to B$ and $g: C \to D$, we obtain a map $[f,g]: [B,C] \to [A,D]$ which is interpreted precisely as pre-composing by $f$ and post-composing by $g$, so $[f,g](-) = f;-;g$. This induces a functor $[-,-]$ which contravariant in its first argument and covariant in its second argument. Now in a setting in which we have a dagger, note that for every map $f: A\to B$, we obtain the map $[f^\dagger, f]: [A,A] \to [B,B]$, which we interpret as conjugation by $f$, so $[f^\dagger, f](-) = f^\dagger; -; f$. Now in a $\dagger$KCC, for every map $f: A \to B$, we get a map $f^\ast: B^\ast \to A^\ast$ \cite[Ex 4.4]{selinger2010survey}. Then we get that $[f,g] = f^\ast \otimes g$, and therefore $[f^\dagger, f] = f_\ast \otimes f$, where recall that $(\_)_\ast = (\_)^{\dagger\ast} = (\_)^{\ast\dagger}$. 

With all this notation, we are finally in a position to review abstract quantum channels in a $\dagger$KCC. 

\begin{definition}\label{Def:UCPTP} In a $\dagger$KCC, a map $\Phi: [A,A] \to [B,B]$ is said to be:
  \begin{enumerate}[(i)]
\item \label{def:CP}\textbf{Completely Positive} \cite[Cor 4.13]{selinger2007dagger} (or \textbf{CP} for short), if there exists a map $f: A \to X \otimes B$ such that $\Phi = [f^\dagger, f]; \mathsf{tr}^X_{B,B}$; 
\item \label{def:TP}\textbf{Trace Preserving} (or \textbf{TP} for short), if $\Phi; \mathsf{tr}^B=\mathsf{tr}^A$; 
\item \label{def:U}\textbf{Unital} (or \textbf{U} for short), if $\mathsf{u}^A;\Phi = \mathsf{u}^B$
\end{enumerate}
A map which is CP and TP will be called a \textbf{CPTP} map, and a map which is CPTP and unital will be called a \textbf{UCPTP} map. 
\end{definition}

\begin{remark} Notice that the definition of being U can be defined in a symmetric monoidal closed category, being TP can be defined in a traced symmetric monoidal closed category, and in theory CP can be defined in a so-called dagger traced symmetric monoidal closed category ($\dagger$TSMCC). As such, an argument could be made to write this story in the setting of a $\dagger$TSMCC. However there a few issues with this at this time. Firstly, it is unclear what exactly should be the proper coherences for a $\dagger$TSMCC. Secondly,   while every $\dagger$KCC is a $\dagger$TSMCC, we do not currently have any good examples of a $\dagger$TSMCC which is not a $\dagger$KCC. This is also a problem for dagger symmetric monoidal closed categories \cite[Sec 9]{heunen2016monads}. We leave the exploration of $\dagger$TSMCC and quantum channels in this setting for future work. 
\end{remark}

\begin{example} In $\mathsf{FHILB}$, the object $[H,H]$ is more commonly denoted by $\mathcal{L}(H)$, the space of all $\mathbb{C}$-linear endomaps on $H$. Then, as expected, the internal trace operator $\mathsf{tr}^H: \mathcal{L}(H) \to \mathbb{C}$ sends an endomap to its trace, $\mathsf{tr}^H(\rho) = \mathsf{Tr}(\rho)$, while the internal partial trace operator $\mathsf{tr}^K_{H,H}:\mathcal{L}(K \otimes H) \to \mathcal{L}(H)$ gives of the partial trace, $\mathsf{tr}^H_{K,K}(\rho) = \mathsf{Tr}_K(\rho)$. On the other hand, $\mathsf{u}^H: \mathbb{C} \to \mathcal{L}(H)$ picks out the identity map, $\mathsf{u}^H(1) = \mathsf{id}_H$. Now recall that a $\mathbb{C}$-linear map $\Phi: \mathcal{L}(H) \to \mathcal{L}(H^\prime)$ is CP in the usual sense if when tensoring with the identity of any dimension, $\mathsf{id} \otimes \Phi$ sends positive semidefinite maps to positive semidefinite maps. Alternatively, one of the well-known equivalent characterization being CP is if there exists an $\psi: H \to K \otimes H^\prime$ such that $\Phi(\rho) = \mathsf{Tr}_K\left(\psi^\dagger \circ \rho \circ \psi\right)$ \cite[Thm 2.22]{watrous2018theory}, which is precisely what Def \ref{Def:UCPTP}.(\ref{def:CP}) corresponds to. Now $\Phi$ is TP if $\mathsf{Tr}(\Phi(\rho)) = \mathsf{Tr}(\rho)$ and U if $\Phi(\mathsf{id}_H) = \mathsf{id}_{H^\prime}$, which correspond precisely to Def \ref{Def:UCPTP}.(\ref{def:TP}) and (\ref{def:U}) respectively. 
\end{example}

\begin{example} In $\mathsf{REL}$, $[X,X]= X \times X$ and so a map $\Phi: [X,X] \to [Y,Y]$ is a subset $\Phi \subseteq X \times X \times Y \times Y$. Then $\Phi$ is CP if it respects inverses \cite[Def 5.2 \& Prop 5.3]{coecke2016categories}, that is, if it satisfies (i) $(x_1,x_2,y_1,y_2) \in R \Leftrightarrow (x_2,x_1,y_2,y_1) \in R$, and (ii) $(x_1,x_2,y_1,y_2) \in R \Rightarrow (x_1,x_1, y_1, y_1) \in R$ \cite[Prop 5.3]{coecke2016categories}. On the other hand, $\Phi$ is TP if it satisfies that for all $x_1, x_2 \in X$, $\exists y \in Y. (x_1, x_2, y, y) \in \Phi \Leftrightarrow x_1 = x_2$, and dually $\Phi$ is U if it satisfies that for all $y_1, y_2 \in Y$, $\exists x \in X. (x,x, y_1, y_2) \Leftrightarrow y_1 = y_2$. 
\end{example}

Now CP maps are a well-studied concept in a $\dagger$KCC and other frameworks. In particular, CP maps are closed under composition, monoidal product, and dagger \cite[Lem 4.17]{selinger2007dagger}. Therefore, we can build the category of CP maps of a $\dagger$KCC, called the CPM construction \cite[Def 4.18]{selinger2007dagger}, and it is itself a $\dagger$KCC \cite[Thm 4.20]{selinger2007dagger}. The CPM construction is arguably one of the most celebrated constructions in categorical quantum foundations, and it has been generalized to other settings \cite{coecke2016categories,coecke2016_cpinf,Coecke_2012,Cunningham_2015,Gogioso_2019}. On the other hand, while TP (resp. U) maps are closed under composition, monoidal product, and the caps and cups are TP (resp. U), the dagger of a TP (resp. U) map is not necessarily TP (resp. U). This is possibly one of the reasons why TP and U have taken less of central role in the study of quantum channels in categorical quantum foundations. However, it is easy to see that since TP and U are dagger adjoint properties, in the sense that we get that the adjoint of a TP is TP if and only if the map was also U to begin with, which is a well known result in quantum information \cite[Thm 2.26]{watrous2018theory}. As such, we can thus conclude that the category of UCPTP maps is $\dagger$KCC (in fact a sub-$\dagger$KCC of the CPM construction). 

\begin{lemma} In a $\dagger$KCC, a map $\Phi: [A,A] \to [B,B]$ is TP (resp. U) if and only if $\Phi^\dagger: [B,B] \to [A,A]$ is U (resp. TP). Therefore, if $\Phi$ is TP (resp. U), then $\Phi^\dagger$ is TP (resp. U) if and only if $\Phi$ is U (resp. TP). 
\end{lemma}
\begin{proof} In our notation, this follows from the fact that $\left({\mathsf{u}^X}\right)^\dagger = \mathsf{tr}^X$ and $\left({\mathsf{tr}^X}\right)^\dagger = \mathsf{u}^X$. 
\end{proof}

\begin{corollary} For a $\dagger$KCC, its category of UCPTP maps is a $\dagger$KCC. 
\end{corollary}

\section{Drazin Inverses}

One of the main result in \cite{cao2021nisq} was that the Drazin inverse of a quantum channel was TP \cite[Thm 2]{cao2021nisq}. The goal of this section is to generalize this result to the setting of $\dagger$KCC. We will show that the Drazin inverse of a TP map is again TP, and that the same result holds for U maps. For an in-depth introduction to Drazin inverses in categories, we refer the reader to \cite{cockett2024drazin}. 

\begin{definition}  \label{def:Drazin}  In a category, a \textbf{Drazin inverse} \cite[Def 2.2]{cockett2024drazin} of an endomap $f: A \to A$ is an endomap $f^D: A \to A$ such that:  
\begin{enumerate}[{\bf [$\text{D}$.1]}]
\begin{multicols}{2}
\item There is a $k \in \mathbb{N}$ such that $f^{k+1} ;f^D = f^k$;
\columnbreak
\begin{multicols}{2} 
\item $f^D ;f ;f^D = f^D$; 
\columnbreak
\item $f ;f^D = f^D ;f$. 
\end{multicols}
\end{multicols}
\end{enumerate}
\vspace{-0.3cm}
If $f$ has a Drazin inverse, we say that $f$ is \textbf{Drazin invertible} or simply \textbf{Drazin}. We call the smallest natural number $k$ such that \textbf{[D.1]} holds the \textbf{Drazin index of $f$}, which we denote by $\mathsf{ind}^D(f)$. A \textbf{Drazin category} is a category such that every map is Drazin. 
\end{definition}

Drazin inverses are \emph{unique} \cite[Prop 2.3]{cockett2024drazin}, so we may speak of \emph{the} Drazin inverse of a map and that being Drazin is a property rather than structure. 
%Many examples and properties of Drazin inverses can be found in \cite{cockett2024drazin}. 

\begin{example} $\mathsf{FHILB}$ is Drazin, so every $\mathbb{C}$-linear map $\rho: H \to H$ admits a Drazin inverse $\rho^D: H \to H$. This follows from the fact that every complex square matrix admits a Drazin inverse \cite[Chap 7]{campbell2009generalized}, where the Drazin inverse computed using the Jordan–Chevalley decomposition \cite[Thm 7.2.1]{campbell2009generalized}. See also \cite[Sec 2.8]{cockett2024drazin} for more details. 
\end{example}

\begin{example} $\mathsf{REL}$ is not Drazin, since in particular the successor relation $\lbrace (n, n+1) \vert~ n \in \mathbb{\mathbb{N}} \rbrace \subseteq \mathbb{N} \times \mathbb{N}$ does not have a Drazin inverse \cite[Sec 2.10]{cockett2024drazin}. However, if we instead consider the subcategory $\mathsf{FREL}$ of \textit{finite sets} and relations, then $\mathsf{FREL}$ is Drazin, since any category whose homsets are finite sets is Drazin \cite[Lem 2.11]{cockett2024drazin}. So if $X$ is a finite set, for every subset $R \subseteq X \times X$ admits a Drazin inverse $R^D \subseteq X \times X$. 
\end{example}

Our goal is now to show that the Drazin inverse, if it exists, of a TP map is again TP. To do so, we make use of the following fact about Drazin inverses. 

\begin{proposition} \label{Drazin-commuting} \cite[Prop 3.10]{cockett2024drazin} In a category, let $f: A \to A$ and $g: B \to B$ be Drazin. For any map $k: A \to B$, if the square on the left commutes, then the square on the right commutes: 
\vspace{-0.3cm}
\[ \begin{array}[c]{c} \begin{tikzcd}
	A & A \\
	B & B
	\arrow["f", from=1-1, to=1-2]
	\arrow["k"', from=1-1, to=2-1]
	\arrow["k", from=1-2, to=2-2]
	\arrow["g"', from=2-1, to=2-2]
\end{tikzcd} \end{array} \Longrightarrow \begin{array}[c]{c} \begin{tikzcd}
	A & A \\
	B & B
	\arrow["f^D", from=1-1, to=1-2]
	\arrow["k"', from=1-1, to=2-1]
	\arrow["k", from=1-2, to=2-2]
	\arrow["g^D"', from=2-1, to=2-2]
\end{tikzcd} \end{array} \]
\end{proposition}

\begin{proposition}\label{prop:DrazinTP} In a $\dagger$KCC, if $\Phi: {[A,A]} \to [A,A]$ is Drazin and also TP, then $\Phi^D: {[A,A]} \to {[A,A]}$ is also TP.
\end{proposition}
\begin{proof} Identity maps $1$ are always Drazin with $1^D = 1$ \cite[Lem 3.6]{cockett2024drazin}. Therefore, by setting $g=1_I$ and $k=\mathsf{tr}^A$ in Prop \ref{Drazin-commuting}, the diagram on the left trivially commutes so the diagram on the right commutes: 
\[\begin{array}[c]{c} \begin{tikzcd}
	{[A,A]} & {[A,A]} \\
	I & I
	\arrow["\Phi", from=1-1, to=1-2]
	\arrow["{\mathsf{tr}^A}"', from=1-1, to=2-1]
	\arrow["{\mathsf{tr}^A}", from=1-2, to=2-2]
	\arrow[from=2-1, to=2-2, equal]
\end{tikzcd} \end{array} \Longrightarrow \begin{array}[c]{c}\begin{tikzcd}
	{[A,A]} & {[A,A]} \\
	I & I
	\arrow["\Phi^D", from=1-1, to=1-2]
	\arrow["{\mathsf{tr}^A}"', from=1-1, to=2-1]
	\arrow["{\mathsf{tr}^A}", from=1-2, to=2-2]
	\arrow[from=2-1, to=2-2, equal]
\end{tikzcd} \end{array} \]
Thus we conclude that the Drazin inverse is TP. 
\end{proof}

So by applying the above result to $\mathsf{FHILB}$, we obtain a novel proof of \cite[Thm 2]{cao2021nisq}, which requires none of the technical linear algebra used in \cite[App A]{cao2021nisq}. Moreover, by a similar argument, we also obtain that the Drazin inverse of U map is again U, which to the best of our knowledge is a new observation. 

\begin{lemma}\label{lemma:DrazinU} In a $\dagger$KCC, if $\Phi: {[A,A]} \to [A,A]$ is U and Drazin, then $\Phi^D: {[A,A]} \to {[A,A]}$ is also U.
\end{lemma}

Therefore, this implies that the Drazin inverse of unital quantum channel is both TP and unital. 

\begin{remark} Notice that that Lemma \ref{lemma:DrazinU} is in fact true in a symmetric monoidal closed category, and similarly Prop \ref{prop:DrazinTP} is also true in a traced symmetric monoidal closed category. 
\end{remark}

In general, however, the Drazin inverse of CP map is not CP. Here is a simple example. 

\begin{example}\label{ex:depolar} Let $\mathbb{R}$ be the field of real numbers and let $H$ be a finite dimensional Hilbert space with dimension $d \geq 1$. Then for any $a \in \mathbb{R}$, define the depolarizing channel $\mathcal{D}_a: \mathcal{L}(H) \to \mathcal{L}(H)$ as $\mathcal{D}_a(\rho) = (1-a)\cdot\rho + \frac{a \mathsf{Tr}(\rho)}{d} \cdot \mathsf{id}_H$. Now for any $a$, $\mathcal{D}_a$ is both TP and U, and if $0 \leq a \leq 1+ \frac{1}{1-d^2}$, then $\mathcal{D}_a$ is also CP -- this is a well known unital quantum channel. Now if $a \neq 0$, then $\mathcal{D}_a$ is an isomorphism with inverse $\mathcal{D}_a^{-1} = \mathcal{D}_{\frac{1}{a}}$. Since the Drazin inverse of an isomorphism is its inverse \cite[Lem 3.6]{cockett2024drazin}, we get $\mathcal{D}^D_a = \mathcal{D}_{\frac{1}{a}}$, which is both TP and U as well. However, the only $a \neq 0$ for which both $a$ and $\frac{1}{a}$ are less than or equal to $1+ \frac{1}{1-d^2}$ is $1$. Therefore, for any $0 < a \leq 1+ \frac{1}{1-d^2}$ and $a \neq 1$, $\mathcal{D}^D_a = \mathcal{D}_{\frac{1}{a}}$ is not CP. 
\end{example}

Now notice that the converse Prop \ref{prop:DrazinTP} and Lemma \ref{lemma:DrazinU} is not necessarily true, in the sense that even if the Drazin inverse is TP or U, the starting map may not be. This because, if $f$ is Drazin, then the Drazin inverse $f^D$ is also Drazin but its Drazin inverse is not necessarily $f$, instead its $f^{DD} = f f^D f$ \cite[Lem 3.11.(i)]{cockett2024drazin}. A Drazin map $f$ is the Drazin inverse of its Drazin inverse $f^D$ precisely when $\mathsf{ind}(f) \leq 1$ \cite[Lem 3.9 \& 3.11.(iii)]{cockett2024drazin}. Explicitly, an endomap $f: A \to A$ is Drazin with $\mathsf{ind}(f) \leq 1$ if and only if there is an endomap $f^D: A \to A$ such that (i) $f ;f^D ;f = f$, (ii) $f^D ;f;f^D = f^D$, and $(iii)$ $f^D ;f = f ;f^D$. In this case, we say that $f^D$ is a \textbf{group inverse} \cite[Def 3.8]{cockett2024drazin} of $f$, and we also have that $f^{DD}=f$. 

\begin{proposition} In a $\dagger$KCC, if $\Phi: {[A,A]} \to [A,A]$ Drazin with $\mathsf{ind}^D(\Phi) \leq 1$ (or equivalently has a group inverse) then $\Phi^D$ is TP (resp. U) if and only if $\Phi^D: {[A,A]} \to {[A,A]}$ is also TP (resp. U).
\end{proposition}
\begin{proof} For the $\Rightarrow$ direction, follows immediately from Prop \ref{prop:DrazinTP} and Lemma \ref{lemma:DrazinU} respectively. While for the $\Leftarrow$, since $\Phi^D$ is Drazin and in this case have that $\Phi^{DD}=\Phi$, then we apply Prop \ref{prop:DrazinTP} and Lemma \ref{lemma:DrazinU} respectively to get that $\Phi$ is TP or U respectively. 
\end{proof}

\section{Dagger Drazin Inverses}

One of the limitations of the Drazin inverse is that it can be computed only for endomaps $A \to A$. Of course, one would like to talk about Drazin inverses of maps of arbitrary type $A \to B$. In \cite{cockett2025dagger}, the authors achieved this by introducing the notion of dagger Drazin inverses for maps in a dagger category. 

\begin{definition}\label{def:dagger-Drazin}  In a dagger category, a \textbf{$\dagger$-Drazin inverse} \cite[Def 2.1]{cockett2025dagger} of a map $f: A \to B$ is a map of dual type $f^\partial: B \to A$ such that:  
\begin{enumerate}[{\bf [$\text{D}^{\dagger}$.1]}]
\item There is a $k \in \mathbb{N}$ such that $(f ;f^\dagger)^k ;f; f^\partial = (f ;f^\dagger)^k$ and $f^\partial ;f ;(f^\dagger; f)^k = (f^\dagger ;f)^k$; 
\vspace{-0.3cm}
\begin{multicols}{3}
\item $f^\partial; f ;f^\partial = f^\partial$; 
\columnbreak
\item $(f;f^\partial)^\dagger = f;f^\partial$;
\columnbreak
\item $(f^\partial ;f)^\dagger = f^\partial; f$.
\end{multicols}
\end{enumerate}
\vspace{-0.3cm}
If $f$ has a $\dagger$-Drazin inverse, we say that $f$ is \textbf{$\dagger$-Drazin}. 
\end{definition}

Now $\dagger$-Drazin inverses are also unique \cite[Prop 2.2]{cockett2025dagger}, so we may speak \textit{the} $\dagger$-Drazin inverse. Moreover, $\dagger$-Drazin inverses are fundamentally linked to Drazin inverses in the following way. 

\begin{theorem}\label{thm:Drazin=dag-Drazin}\cite[Thm 3.6]{cockett2025dagger}  In a dagger category, for a map $f$, the following are equivalent: 
\begin{enumerate}[(i)]
\begin{multicols}{4}
\item $f$ is $\dagger$-Drazin;
\columnbreak
\item $f^\dagger$ is $\dagger$-Drazin;
\columnbreak
\item $f ;f^\dagger$ is Drazin; 
\columnbreak
\item $f^\dagger ;f$ is Drazin.
\end{multicols}
\end{enumerate}
Moreover, if $f$ is $\dagger$-Drazin, then:
\vspace{-0.3cm}
\begin{align*}
    f^\partial = f^\dagger ;(f ;f^\dagger)^D = (f^\dagger ;f)^D ;f^\dagger && {f^\dagger}^\partial = f ;(f^\dagger ;f)^D  =  (f ;f^\dagger)^D ;f && (f ;f^\dagger)^D = {f^\dagger}^\partial; f^\partial && (f^\dagger ;f)^D = f^\partial ;{f^\dagger}^\partial
\end{align*}
\end{theorem}

As such, the above characterization of $\dagger$-Drazin inverses implies that in a dagger category which is also Drazin, every map has a $\dagger$-Drazin inverse. 
%Thus every map in $\mathsf{FHILB}$ and $\mathsf{FREL}$ has a $\dagger$-Drazin inverse. Moreover, as we will review in the next section, the Moore-Penrose inverse is a special case of the $\dagger$-Drazin inverse. 

Now, as expected, the $\dagger$-Drazin inverse of a CP map need not be CP, Ex \ref{ex:depolar} provides an example of such (since in this case, the Drazin inverse and $\dagger$-Drazin inverse coincide). On the other hand, we now wish to generalize that Drazin inverse preserves being TP and U for $\dagger$-Drazin inverses. However, we quickly come across an issue since we know that the Moore-Penrose inverse (which as we will soon review is a special case of a $\dagger$-Drazin inverse) of a TP map need not be TP \cite[Ex 2]{cao2021nisq}. To resolve this issue, as we will see below, it turns that if a map is both TP and U, then its $\dagger$-Drazin inverse will be TP and U. In order to prove this, we first generalize \cite[Prop 3.10]{cockett2024drazin} for $\dagger$-Drazin inverses. 
%If we have space, give the example. 

\begin{proposition}\label{dagger-Drazin-commuting} In a dagger category, suppose that $f: A \to B$ and $g: C \to D$ are $\dagger$-Drazin. Then if the two diagrams on the left commute, then the two diagrams on the right commute: 
\vspace{-0.3cm}
\[ \begin{array}[c]{c} \begin{tikzcd}
	A & B & B & A \\
	C & D & D & C
    	\arrow["{(\textbf{a})}"{description}, draw=none, from=1-1, to=2-2]
        \arrow["{(\textbf{b})}"{description}, draw=none, from=1-3, to=2-4]
	\arrow["f", from=1-1, to=1-2]
	\arrow["h"', from=1-1, to=2-1]
	\arrow["k", from=1-2, to=2-2]
	\arrow["{f^\dagger}", from=1-3, to=1-4]
	\arrow["k"', from=1-3, to=2-3]
	\arrow["h", from=1-4, to=2-4]
	\arrow["g"', from=2-1, to=2-2]
	\arrow["{g^\dagger}"', from=2-3, to=2-4]
\end{tikzcd} \end{array} \Longrightarrow \begin{array}[c]{c} \begin{tikzcd}
	B & A & A & B \\
	D & C & C & D
    \arrow["{(\textbf{c})}"{description}, draw=none, from=1-1, to=2-2]
        \arrow["{(\textbf{d})}"{description}, draw=none, from=1-3, to=2-4]
	\arrow["{f^\partial}", from=1-1, to=1-2]
	\arrow["k"', from=1-1, to=2-1]
	\arrow["h", from=1-2, to=2-2]
	\arrow["{f^{\partial\dagger}= f^{\dagger\partial}}", from=1-3, to=1-4]
	\arrow["h"', from=1-3, to=2-3]
	\arrow["k", from=1-4, to=2-4]
	\arrow["{g^\partial}"', from=2-1, to=2-2]
	\arrow["{g^{\partial\dagger}= g^{\dagger\partial}}"', from=2-3, to=2-4]
\end{tikzcd} \end{array} \]
\vspace{-0.5cm}
\end{proposition}
\begin{proof} By pasting (\textbf{a}) and (\textbf{b}) together, we obtain the diagrams (\textbf{e}) and (\textbf{f}) below. By Thm \ref{thm:Drazin=dag-Drazin}, we know that $f;f^\dagger$, $g;g^\dagger$, $f^\dagger;f$, and $g^\dagger;g$ are all Drazin. Then by Prop \ref{Drazin-commuting}, since (\textbf{e}) and (\textbf{f}) commute, we get that (\textbf{g}) and (\textbf{h}) commute respectively: 
\[ \begin{array}[c]{c} \begin{tikzcd}
	A & A \\
	C & C
    \arrow["{(\textbf{e})}"{description}, draw=none, from=1-1, to=2-2]
	\arrow["f;f^\dagger", from=1-1, to=1-2]
	\arrow["h"', from=1-1, to=2-1]
	\arrow["h", from=1-2, to=2-2]
	\arrow["g;g^\dagger"', from=2-1, to=2-2]
\end{tikzcd} \end{array} \Longrightarrow \begin{array}[c]{c} \begin{tikzcd}
	A & A \\
	C & C
    \arrow["{(\textbf{g})}"{description}, draw=none, from=1-1, to=2-2]
	\arrow["(f;f^\dagger)^D", from=1-1, to=1-2]
	\arrow["h"', from=1-1, to=2-1]
	\arrow["h", from=1-2, to=2-2]
	\arrow["(g;g^\dagger)^D"', from=2-1, to=2-2]
\end{tikzcd} \end{array}  ,\qquad \begin{array}[c]{c} \begin{tikzcd}
	B & B \\
	D & D
    \arrow["{(\textbf{f})}"{description}, draw=none, from=1-1, to=2-2]
	\arrow["f^\dagger;f", from=1-1, to=1-2]
	\arrow["k"', from=1-1, to=2-1]
	\arrow["k", from=1-2, to=2-2]
	\arrow["g^\dagger;g"', from=2-1, to=2-2]
\end{tikzcd} \end{array} \Longrightarrow \begin{array}[c]{c} \begin{tikzcd}
	B & B \\
	D & D
         \arrow["{(\textbf{h})}"{description}, draw=none, from=1-1, to=2-2]
	\arrow["(f^\dagger;f)^D", from=1-1, to=1-2]
	\arrow["k"', from=1-1, to=2-1]
	\arrow["k", from=1-2, to=2-2]
	\arrow["(g^\dagger;g)^D"', from=2-1, to=2-2]
\end{tikzcd} \end{array}
\]
Then by pasting (\textbf{g}) with (\textbf{b)} and (\textbf{h}) with (\textbf{a}), we get that the following diagrams commute: 
\[
  \begin{tikzcd}
	A & A & B \\
	C & C & D
	\arrow["{(f;f^\dagger)^D}", from=1-1, to=1-2]
	\arrow["h"', from=1-1, to=2-1]
	\arrow["{{(\textbf{g})}}"{description}, draw=none, from=1-1, to=2-2]
	\arrow["f", from=1-2, to=1-3]
	\arrow["h", from=1-2, to=2-2]
	\arrow["{{(\textbf{a})}}"{description}, draw=none, from=1-2, to=2-3]
	\arrow["k", from=1-3, to=2-3]
	\arrow["{(g;g^\dagger)^D}"', from=2-1, to=2-2]
	\arrow["g"', from=2-2, to=2-3]
\end{tikzcd} \qquad  \begin{tikzcd}
	B & B & A \\
	D & D & B
	\arrow["{(f^\dagger;f)^D}", from=1-1, to=1-2]
	\arrow["k"', from=1-1, to=2-1]
	\arrow["{{(\textbf{h})}}"{description}, draw=none, from=1-1, to=2-2]
	\arrow["{f^\dagger}", from=1-2, to=1-3]
	\arrow["k", from=1-2, to=2-2]
	\arrow["{{(\textbf{b})}}"{description}, draw=none, from=1-2, to=2-3]
	\arrow["h", from=1-3, to=2-3]
	\arrow["{(g^\dagger;g)^D}"', from=2-1, to=2-2]
	\arrow["{g^\dagger}"', from=2-2, to=2-3]
\end{tikzcd} \]
However again by Thm \ref{thm:Drazin=dag-Drazin} we have that $f^\partial = (f^\dagger ;f)^D ;f^\dagger$, $g^\partial = (g^\dagger ;g)^D ;g^\dagger$, ${f^\dagger}^\partial = (f ;f^\dagger)^D ;f$, and ${g^\dagger}^\partial = (g ;g^\dagger)^D ;g$, these two diagrams above are (\textbf{d}) and (\textbf{c}) respectively. 
\end{proof}

\begin{proposition}\label{dag-Drazin-TP} In a $\dagger$KCC, if $\Phi: {[A,A]} \to {[B,B]}$ is $\dagger$-Drazin, and both TP and U, then $\Phi^\partial: {[B,B]} \to {[A,A]}$ is TP and U. 
\end{proposition}
\begin{proof} Identity maps $1$ are always $\dagger$-Drazin with $1^\partial = 1$ \cite[Lem 2.5]{cockett2025dagger}. Therefore, by setting $g=1_I$, $h=\mathsf{tr}^A$, and $k=\mathsf{tr}^B$ in Prop \ref{dagger-Drazin-commuting}, we get that the two diagrams on the left commute trivially, which imply that the two diagrams on the right commute. 
    \[ \begin{array}[c]{c} \begin{tikzcd}
	{[A,A]} & {[B,B]} & {[B,B]} & {[A,A]} \\
	I & I & I & I
	\arrow["\Phi", from=1-1, to=1-2]
	\arrow["{\mathsf{tr}^A}"', from=1-1, to=2-1]
	\arrow["{\mathsf{tr}^B}", from=1-2, to=2-2]
	\arrow["{{\Phi^\dagger}}", from=1-3, to=1-4]
	\arrow["{\mathsf{tr}^B}"', from=1-3, to=2-3]
	\arrow["{\mathsf{tr}^A}", from=1-4, to=2-4]
	\arrow[from=2-1, to=2-2, equal]
	\arrow[from=2-3, to=2-4, equal]
\end{tikzcd} \end{array} \Longrightarrow \begin{array}[c]{c} \begin{tikzcd}
	{[B,B]} & {[A,A]} & {[A,A]} & {[B,B]} \\
	I & I & I & I
	\arrow["{\Phi^\partial}", from=1-1, to=1-2]
	\arrow["{\mathsf{tr}^B}"', from=1-1, to=2-1]
	\arrow["{\mathsf{tr}^A}", from=1-2, to=2-2]
	\arrow["{\Phi^{\dagger\partial}= \Phi^{\partial\dagger}}", from=1-3, to=1-4]
	\arrow["{\mathsf{tr}^B}"', from=1-3, to=2-3]
	\arrow["{\mathsf{tr}^A}", from=1-4, to=2-4]
	\arrow[from=2-1, to=2-2, equal]
	\arrow[from=2-3, to=2-4, equal]
\end{tikzcd} \end{array} \]
So we conclude that the $\dagger$-Drazin $\Phi^\partial$ is TP and U as desired. 
\end{proof}

We note that the converse of the proposition is not true. The reason is similar as to why for Drazin inverses the converse of Prop \ref{prop:DrazinTP} wasn't necessarily true. Indeed, while if $f$ is a $\dagger$-Drazin, its $\dagger$-Drazin inverse $f^\partial$ is also $\dagger$-Drazin, $f^{\partial \partial}$ is not necessarily $f$, instead its $f^{\partial \partial}= f;f^\partial;f$ \cite[Prop 2.6.(iv)]{cockett2025dagger}. This issue is solved when we instead consider the \textit{Moore-Penrose inverse}.  

\section{Moore-Penrose Inverses}

Arguably, the most well-known example of a generalized inverse is the famous Moore-Penrose inverse. However, as explained in \cite{cao2021nisq}, an issue with the Moore-Penrose inverse is that, unlike the Drazin inverse, the Moore-Penrose inverse of a TP map need not be TP, see \cite[Ex 2]{cao2021nisq}. In this section, we explain how to fix this issue: we need to assume U. In fact, we will show that a map that has a Moore-Penrose inverse is TP and U if and only if its Moore-Penrose inverse is TP and U. For a detailed introduction to Moore-Penrose inverses in dagger categories, we refer the reader to see \cite{Cockett2023Moore}.

\begin{definition} In a dagger category, a \textbf{Moore-Penrose inverse} \cite[Def 2.3]{Cockett2023Moore} of a map $f: A \to B$ is a map of dual type $f^\circ: B \to A$ such that:
\vspace{-0.3cm}
\begin{enumerate}[{\bf [$\text{MP}$.1]}]
\begin{multicols}{4}
    \item $f f^\circ f = f$
\columnbreak
    \item $f^\circ f f^\circ = f^\circ$
\columnbreak
\item $(ff^\circ)^\dagger = ff^\circ$
\columnbreak
    \item $(f^\circ f)^\dagger = f^\circ f$
\end{multicols}
\end{enumerate}
\vspace{-0.3cm}
If $f$ has a Moore-Penrose inverse, we say that $f$ is \textbf{Moore-Penrose} (or \textbf{MP} for short). A \textbf{Moore-Penrose dagger category} is a dagger category such that every map is MP. 
\end{definition}

Moore-Penrose inverses are \emph{unique} \cite[Lem 2.4]{Cockett2023Moore}, so again we may speak of \emph{the} Moore-Penrose inverse of a map. Moreover, if $f$ is MP, then its MP inverse $f^\circ$ is also MP where $f^{\circ\circ} = f$ \cite[Lem 2.5.(i)]{Cockett2023Moore} (which, as we have seen, is not necessarily true for other kinds of generalzied inverses). 

%Many examples and properties of Moore-Penrose inverses and Moore-Penrose dagger categories can be found in \cite{Cockett2023Moore}. 

\begin{example} $\mathsf{FHILB}$ is Moore-Penrose, so every $\mathbb{C}$-linear map $f: H \to K$ admits a MP inverse $f^\circ: K \to H$ which can be constructed using singular value decomposition \cite[Ex 2.13]{Cockett2023Moore}. 
\end{example}

\begin{example} $\mathsf{FREL}$ is not Moore-Penrose but we do know which maps are MP. Recall that a relation $R \subseteq X \times Y$ is \textit{difunctional} if $(x,b),(a,b),(a,y) \in R$ implies that $(x,y) \in R$ (which from a dagger category perspective means that $R$ is a partial isometry). Then a relation $R \subseteq X \times Y$ is MP if and only if it difunctional, and in this case the MP inverse is the inverse relation, $R^\circ = R^\dagger$ \cite[Ex 2.14]{Cockett2023Moore}. 
\end{example}

It turns out that an MP inverse is a special case of the $\dagger$-Drazin inverse. In fact, MP maps are precisely the $\dagger$-Drazin maps that are the $\dagger$-Drazin inverse of their $\dagger$-Drazin inverse. 

\begin{theorem}\label{thm:MP=a-dag-Drazin} \cite[Thm 4.2]{cockett2025dagger} In a dagger category, a map $f$ is MP if and only if $f$ is $\dagger$-Drazin and $f^{\partial \partial} = f$. In this case, the MP inverse and $\dagger$-Drazin inverse coincide, $f^\partial = f^\circ$. 
\end{theorem}

With this in hand, we can easily extend Prop \ref{dagger-Drazin-commuting} to MP inverses, and show that MP inverses preserve being both TP and U, as desired. 

\begin{proposition}\label{MP-commuting} Suppose that $f: A \to B$ and $g: C \to D$ are MP. Then, if the two diagrams on the left commute, then the two diagrams on the right commute: 
\vspace{-0.3cm}
\[ \begin{array}[c]{c} \begin{tikzcd}
	A & B & B & A \\
	C & D & D & C
	\arrow["f", from=1-1, to=1-2]
	\arrow["h"', from=1-1, to=2-1]
	\arrow["k", from=1-2, to=2-2]
	\arrow["{f^\dagger}", from=1-3, to=1-4]
	\arrow["k"', from=1-3, to=2-3]
	\arrow["h", from=1-4, to=2-4]
	\arrow["g"', from=2-1, to=2-2]
	\arrow["{g^\dagger}"', from=2-3, to=2-4]
\end{tikzcd} \end{array} \Longrightarrow \begin{array}[c]{c} \begin{tikzcd}
	B & A & A & B \\
	D & C & C & D
	\arrow["{f^\circ}", from=1-1, to=1-2]
	\arrow["k"', from=1-1, to=2-1]
	\arrow["h", from=1-2, to=2-2]
	\arrow["{f^{\circ\dagger}= f^{\dagger\circ}}", from=1-3, to=1-4]
	\arrow["h"', from=1-3, to=2-3]
	\arrow["k", from=1-4, to=2-4]
	\arrow["{g^\circ}"', from=2-1, to=2-2]
	\arrow["{g^{\circ\dagger}= g^{\dagger\circ}}"', from=2-3, to=2-4]
\end{tikzcd} \end{array} \]
\end{proposition}
\begin{proof} Since $f$ and $g$ are MP, they are both $\dagger$-Drazin with $f^\partial= f^\circ$ and $g^\partial=g^\circ$. Then by applying Prop \ref{dagger-Drazin-commuting}, we obtain the desired result. 
\end{proof}

\begin{proposition}\label{MP-TP} In a $\dagger$KCC, if $\Phi: {[A,A]} \to {[B,B]}$ is MP, then $\Phi$ is TP and U if and only if $\Phi^\circ: {[B,B]} \to {[A,A]}$ is TP and U. 
\end{proposition}
\begin{proof} For the $\Rightarrow$ direction, since $\Phi$ is MP, it is also $\dagger$-Drazin and $\Phi^\partial=\Phi^\circ$. By applying Prop \ref{dag-Drazin-TP} we get that $\Phi^\circ$ is TP and U. For the $\Leftarrow$ direction, since $\Phi$ is MP, $\Phi^{\circ}$ is also MP and $\Phi^{\circ\circ} = \Phi$. So by applying the $\Rightarrow$ direction that we just proved with $\Phi^\circ$, we get $\Phi$ is TP and U as desired. 
\end{proof}

Therefore this tells us that for unital quantum channels, the MP inverse is both TP and unital. 

\section{(U)CPTP Maps whose Generalized Inverses are (U)CPTP}

Now, as we have established, while generalized inverses preserve being TP and/or U, they often do not preserve being CP. Thus, a natural question to ask is whether there are any quantum channels whose generalized inverse are again quantum channels. The goal of this section is to provide some examples of (U)CPTP maps that are again (U)CPTP. Our first source of examples are the so called pure channels
%, which recall amount to conjugating operators. 

\begin{definition} In a $\dagger$KCC, a map $\Phi: [A,A] \to [B,B]$ is \textbf{pure} if there exists a map $f: A \to B$ such that $\Phi = [f^\dagger,f]$. 
%We say that $\Phi$ is \textbf{pure unitary} if there exists a unitary map $u: A \to B$ such that $\Phi = [u^\dagger,u]$. 
\end{definition}

The properties of $f$ determine the properties of the induced map $[f^\dagger,f]$. Recall that in a dagger category, a map $f: A \to B$ is an \textbf{isometry} if $f;f^\dagger = 1_A$, a \textbf{coisometry} if $f;f^\dagger= 1_B$, and \textbf{unitary} if it both an isometry and a coisometry, which particular means it is an isomorphism whose inverse is its adjoint, $f^{-1} = f^\dagger$. Similarly, we say that map $\Phi: [A,A] \to [B,B]$ is \textbf{pure isometry (resp. unitary)} if there exists a isometry (resp. unitary) map $f: A \to B$ such that $\Phi = [f^\dagger,f]$. 

\begin{lemma} In a $\dagger$KCC, for any map $f: A \to B$, 
\vspace{-0.3cm}
   \begin{enumerate}[(i)]
   \begin{multicols}{2}
\item $[f^\dagger,f]$ is CP \cite[Lem 4.17.(d)]{selinger2007dagger};
\item  $[f^\dagger,f]$ is TP if and only if $f$ is an isometry;
\columnbreak
\item $[f^\dagger,f]$ is U if and only if $f$ is a coisometry;
\item $[f^\dagger,f]$ is UCPTP if and only if $f$ is unitary. 
\end{multicols}
\end{enumerate}
\vspace{-0.3cm}
Therefore, a pure isometry map is CPTP and a pure unitary map is UCPTP. 
\end{lemma}
\begin{proof} (i) is immediate by construction. For (ii)-(iv), these follow from the fact that since $\mathsf{tr}^X$ and $\mathsf{u}^X$ are extranatural, we have that $[f^\dagger;f];\mathsf{tr}^B= [1_A, f;f^\dagger];\mathsf{tr}^A = [f^\dagger;f, 1_B];\mathsf{tr}^B$ and $\mathsf{u}^A;[f^\dagger;f]= \mathsf{u}^A;[1_A, f;f^\dagger] = \mathsf{u}^B;[f^\dagger;f, 1_B]$. 
\end{proof}

\begin{lemma} In a $\dagger$KCC, if a map $f$ is Drazin or $\dagger$-Drazin or MP, then $[f^\dagger, f]$ is Drazin where $[f^\dagger, f]^D = [{f^D}^\dagger, f^D]$ or $\dagger$-Drazin where $[f^\dagger, f]^\partial = [{f^\partial}^\dagger, f^\partial]$ or MP where $[f^\dagger, f]^\circ= [{f^\circ}^\dagger, f^\circ]$ respectively. In particular, if $f$ is unitary, then $[f^\dagger, f]$ is Drazin, $\dagger$-Drazin, and MP where $[f^\dagger, f]^D = [f^\dagger, f]^\partial = [f^\dagger, f]^\circ= [f, f^\dagger]$. 
\end{lemma}
\begin{proof} This easily follows from the fact that Drazin inverses are absolute, meaning any functors preserves Drazin inverses. It is easy to see that $\dagger$-Drazin inverses or MP inverses are dagger absolute, that is, any dagger functor preserves $\dagger$-Drazin inverses or MP inverses.
\end{proof}

A well-known result about quantum channels is that they admit a Kraus representation \cite[Cor 2.27]{watrous2018theory}, meaning that they are all sums of pure isometric channels. To consider these in our setting, we will have to add additive structure to our $\dagger$KCC. So for a category with biproducts \cite[Sec 6.3]{selinger2010survey}, we denote the zero object by $\mathsf{0}$, and the biproduct of a family of objects by $A_1 \oplus \hdots A_n$ with projection maps $\pi_j: A_1 \oplus \hdots A_n \to A_j$ and injection maps $\iota_j: A_j \to A_1 \oplus \hdots A_n$. Recall that every category with biproducts is also enriched over commutative monoids, which essentially means we can add parallel maps $f+g$ and have zero maps $0$, such that composition is a monoid morphism. Now a dagger category has dagger biproducts \cite[Sec 7.6]{selinger2010survey} if it has biproducts such that the projection and injections are adjoints of each other, so $\pi_i^\dagger=\iota_i$ and $\iota^\dagger_i = \pi_i$. In this setting, the dagger is also a monoid morphism. Then by a \textbf{dagger additive compact closed category} (or \textbf{$\dagger$AKCC} for short), we mean a $\dagger$KCC with dagger biproducts. It is worth noting that in this setting, it comes for free that monoidal product preserves both the sum and the biproduct in the expected way. 

%\begin{example} $\mathsf{FHILB}$ is a $\dagger$AKCC where the biproduct is given by the direct sum of Hilbert spaces, $H \oplus K$, which in turn induces the usual sum of $\mathbb{C}$-linear maps. \end{example}

%\begin{example} $\mathsf{REL}$ is a $\dagger$AKCC where the biproduct is given by the disjoint unit of sets, $X \sqcup Y$, while the induced sum is given by the union, $R \cup S$. \end{example}


\begin{definition} In a $\dagger$AKCC, a map $\Phi: [A,A] \to [B,B]$ is said to be \textbf{mixed pure} if there exits a family of map $\lbrace f_i: A \to B \rbrace^n_{i=1}$ such that $\Phi = \sum\limits_{i=1}^n [f_i^\dagger,f_i]$. 
%In this case, we call $\lbrace f_i: A \to B \rbrace^n_{i=1}$ a family of \textbf{Kraus maps} of $\Phi$. 
\end{definition}

\begin{lemma} In a $\dagger$AKCC, let $\lbrace f_i: A \to B \rbrace^n_{i=1}$ a family of maps. Then: 
\vspace{-0.3cm}
   \begin{enumerate}[(i)]
    \begin{multicols}{2}
\item $\sum\limits_{i=1}^n [f_i^\dagger,f_i]$ is CP;
\item $\sum\limits_{i=1}^n [f_i^\dagger,f_i]$ is TP if and only if $\sum\limits_{i=1}^n f_i;f_i^\dagger = 1_A$;
\columnbreak
\item $\sum\limits_{i=1}^n [f_i^\dagger,f_i]$ is U if and only if $\sum\limits_{i=1}^n f_i^\dagger;f_i = 1_B$. 
\end{multicols}
\end{enumerate}
\vspace{-0.3cm}
\end{lemma}
\begin{proof} For (i), this follows from the fact that sum of CP maps is again a CP map \cite[Lem 5.2]{selinger2007dagger}. The other two statements are easy to check.  
\end{proof}

Now as with the usual inverse, computing the generalized inverse of a sum can be quite complicated \cite{cline1965representations}, and in general the generalized inverse of the sum need not be the sum of the generalized inverses. However, if we assume minimal orthogonal condition on our family of maps, then we can show that in this case, the generalized inverse of their sum is the sum of their generalized inverse. Thus we get the following result in a commutative monoid enriched category. 

\begin{lemma}\label{Drazin-sum} In a category enriched over commutative monoids, let $\lbrace f_i: A \to A \rbrace^n_{i=1}$ be a family of Drazin endomaps such that $f_i f_j=0$ for all $i \neq j$. Then $\sum \limits^n_{i=1} f_i$ is Drazin where $\left(\sum \limits^n_{i=1} f_i \right)^D = \sum \limits^n_{i=1} f_i^D$.
\end{lemma}
\begin{proof} We need to show that $\sum \limits^n_{i=1} f_i^D$ satisfies \textbf{[D.1]-[D.3]}. Let us start with \textbf{[D.3]}. Since $f_i;f_j = 0 = f_j;f_i$, by Prop \ref{Drazin-commuting} and that $0$ is Drazin with $0^D=0$ \cite[Lem 6.10]{cockett2024drazin}, it follows that $f_i^D;f_j = 0 = f_j^D;f_i$. So we get that $\left(\sum \limits^n_{i=1} f_i \right) ; \left(\sum \limits^n_{i=1} f_i^D \right) = \sum \limits^n_{i=1} f_i;f_i^D = \sum \limits^n_{i=1} f_i^D ;f_i = \left(\sum \limits^n_{i=1} f_i \right) ; \left(\sum \limits^n_{i=1} f_i^D \right)$. So \textbf{[D.3]} holds. Moreover, note that $f_j; f_i; f_i^D =0$, then it follows that $\left(\sum \limits^n_{i=1} f_i^D \right);\left(\sum \limits^n_{i=1} f_i \right) ; \left(\sum \limits^n_{i=1} f_i^D \right) = \sum \limits^n_{i=1} f_i^D ;f_i;f_i^D =\sum \limits^n_{i=1} f_i^D$. So \textbf{[D.2]} holds. Now for \textbf{[D.1]}, first note that in our case we have that $\left(\sum \limits^n_{i=1} f_i \right)^{k+1} = \sum\limits^{n}_{i=1} f_i^{k+1}$. Then set $k=\max\lbrace \mathsf{ind}^D(f_i) \rbrace$. In particular this means that $f_i^{k+1} f^D_i = f^{k}_i$ \cite[Lem 2.3.(i)]{cockett2024drazin}. Therefore, since $f_i^{k+1};f_j^D =0$, we have that $\left(\sum \limits^n_{i=1} f_i \right)^{k+2}; \left(\sum \limits^n_{i=1} f_i^D \right) = \sum \limits_{i=1}^n f^{k+2}_i;f_i^D = \sum \limits_{i=1}^n f^{k+1}_i =\left(\sum \limits^n_{i=1} f_i \right)^{k+1}$. So \textbf{[D.1]} holds. Therefore $\left(\sum \limits^n_{i=1} f_i \right)^D = \sum \limits^n_{i=1} f_i^D$ is the Drazin inverse of $\sum \limits^n_{i=1} f_i$ as desired. 
\end{proof}

With this in hand, we can build a (U)CPTP map whose Drazin inverse is also (U)CPTP. 

\begin{corollary}\label{cor:draz-phi-sum} In a $\dagger$AKCC, let $\lbrace f_i: A \to A \rbrace^n_{i=1}$ be a family of Drazin endomaps such that $f_i f_j=0$ for all $i \neq j$, and let $\Phi = \sum\limits_{i=1}^n [f_i^\dagger,f_i]$. Then $\Phi$ is CP and Drazin with $\Phi^D = \sum\limits_{i=1}^n [f_i^{D\dagger},f_i]$, and thus $\Phi^D$ is also CP. Furthermore, if $\sum\limits_{i=1}^n f_i;f_i^\dagger = 1_A$ (resp. $\sum\limits_{i=1}^n f_i^\dagger;f_i = 1_A$), then both $\Phi$ and $\Phi^D$ are TP (resp. U). 
\end{corollary}

Let us now extend Lem \ref{Drazin-sum} to $\dagger$-Drazin inverses and to MP inverse. 

\begin{lemma}\label{lem:dag-draz-sum} In a dagger category enriched over commutative monoids, let $\lbrace f_i: A \to B \rbrace^n_{i=1}$ be a family of $\dagger$-Drazin maps such that $f_i; f^\dagger_j=0$ for all $i \neq j$. Then $\sum \limits^n_{i=1} f_i$ is $\dagger$-Drazin where $\left(\sum \limits^n_{i=1} f_i \right)^\partial = \sum \limits^n_{i=1} f_i^\partial$.
\end{lemma}
\begin{proof} First observe that since $f_i; f^\dagger_j=0$ for all $i \neq j$, we get that $\left( \sum \limits^n_{i=1} f_i\right); \left( \sum \limits^n_{i=1} f_i\right)^\dagger= \sum \limits^n_{i=1} f_i;f_i^\dagger$. Now since each $f_i$ is $\dagger$-Drazin, we have that $f_i;f_i^\dagger$ is Drazin by Thm \ref{thm:Drazin=dag-Drazin}. Moreover, we have that since $f_i; f^\dagger_j=0$ for all $i \neq j$, we also get that $f_i;f_i^\dagger;f_j;f_j^\dagger=0$ for all $i \neq j$. Thus by Lem \ref{Drazin-sum}, we get that $\left( \left( \sum \limits^n_{i=1} f_i\right); \left( \sum \limits^n_{i=1} f_i\right)^\dagger \right)^D = \left(\sum \limits^n_{i=1} f_i;f_i^\dagger\right)^D = \sum \limits^n_{i=1} (f_i;f_i^\dagger)^D$. Therefore we know that $\sum \limits^n_{i=1} f_i$ is also $\dagger$-Drazin via the formula from Thm \ref{thm:Drazin=dag-Drazin}, which is easily worked out to be $\left(\sum \limits^n_{i=1} f_i \right)^\partial = \sum \limits^n_{i=1} f_i^\partial$ as desired. 
%Before we do so, first note that since $f_j^\dagger; f_i; f_i^\dagger = 0 = 0;f_j^\dagger$, and $0$ is clearly always Drazin with $0^D=0$, it follows that $f_j^\dagger;(f_i;f_i^\dagger)^D =0$. Therefore we can compute that: 
%\[ \left(\sum \limits^n_{i=1} f_i \right)^\partial = \left( \sum \limits^n_{i=1} f_i\right)^\dagger\left( \left( \sum \limits^n_{i=1} f_i\right); \left( \sum \limits^n_{i=1} f_i\right)^\dagger \right)^D = \left( \sum \limits^n_{i=1} f_i\right)^\dagger \sum \limits^n_{i=1} (f_i;f_i^\dagger)^D = \sum \limits^n_{i=1} f^\dagger_i;(f_i;f_i^\dagger)^D = \sum \limits^n_{i=1} f_i^\partial\]
%as desired. 
\end{proof}

\begin{corollary} In a dagger category enriched over commutative monoids, let $\lbrace f_i: A \to B \rbrace^n_{i=1}$ be a family of MP maps such that $f_i; f^\dagger_j=0$ for all $i \neq j$. Then $\sum \limits^n_{i=1} f_i$ is MP where $\left(\sum \limits^n_{i=1} f_i \right)^\circ = \sum \limits^n_{i=1} f_i^\circ$.
\end{corollary} 
\begin{proof} This follows immediately from the fact that MP inverses are special cases of $\dagger$-Drazin inverses and applying Lem \ref{lem:dag-draz-sum}.   
\end{proof}

\begin{corollary}\label{cor:dag-draz-phi-sum} In a $\dagger$AKCC, let $\lbrace f_i: A \to B \rbrace^n_{i=1}$ be a family of $\dagger$-Drazin (resp. MP) maps such that $f_i f_j=0$ for all $i \neq j$, and let $\Phi = \sum\limits_{i=1}^n [f_i^\dagger,f_i]$. Then $\Phi$ is CP and $\dagger$-Drazin (resp. MP), and also $\Phi^\partial = \sum\limits_{i=1}^n [f_i^{\partial\dagger},f_i]$ (resp. $\Phi^\circ = \sum\limits_{i=1}^n [f_i^{\circ\dagger},f_i]$), and thus $\Phi^\partial$ (resp. $\Phi^\circ$) is CP. Furthermore, if $\sum\limits_{i=1}^n f_i;f_i^\dagger = 1_A$ (resp. $\sum\limits_{i=1}^n f_i^\dagger;f_i = 1_B$), then both $\Phi$ and $\Phi^\partial$ (resp. $\Phi^\circ$) are TP (resp. U). 
\end{corollary}

It is easy enough to find examples that satisfy the necessary conditions in Cor \ref{cor:draz-phi-sum} and Cor \ref{cor:dag-draz-phi-sum}. Indeed, given a biproduct $\bigoplus\limits^n_{i=1} A_i$, consider the maps $e_i=\pi_i;\iota_i$. Each $e_i$ is self-adjoint and idempotent, and moreover, since we have dagger biproducts we have that $e_i;e_j=0$ when $i \neq j$ and also $\sum^n_{i=1} e_i^2 = 1_{\bigoplus\limits^n_{i=1} A_i}$. Therefore $\sum\limits^n_{i=1} \left[ e_i^\dagger, e_i \right]: \left[\bigoplus\limits^n_{i=1} A_i, \bigoplus\limits^n_{i=1} A_i\right] \to \left[\bigoplus\limits^n_{i=1} A_i, \bigoplus\limits^n_{i=1} A_i\right]$ is a UCPTP, which turns out to be its own Drazin/$\dagger$-Drazin/MP inverse. So this gives us an example of an abstract quantum channel, whose generalized inverse is a quantum channel. 

% \section{Conclusion}

% Conclusion goes here. I don't think we need one. Only if we have space for a short paragraph. 

\newpage 
\nocite{*}
\bibliographystyle{eptcs}
\bibliography{qpl2026}
\end{document}

\note{What is positive definite in the following definition?}
\begin{definition} A \textbf{partition of unity} is a finite family of scalars $\lbrace a_i: I \to I \rbrace^n_{i=1}$ such each $a_i$ is positive definite and $\sum \limits_i a_i = 1_I$. 
\end{definition}

\begin{definition} A map $\Phi: [A,A] \to [B,B]$ is said to be \textbf{mixed unitary} if there exits a partition of unity $\lbrace a_i: I \to I \rbrace^n_{i=1}$ and a family of unitary maps $\lbrace u_i: A \to B \rbrace^n_{i=1}$ such that $\Phi = \sum\limits_{i=1}^n a_i \cdot [u^\dagger_i, u_i]$. 
\end{definition}

\begin{proposition} A mixed unitary map is UCPTP. 
\end{proposition}

\subsection{Depolarizing Maps}

\note{The quantum textbooks say depolarizing maps are mixed unitary! We need to look into this.}

In a traced symmetric monoidal category, the \textbf{dimension} of an object $A$ is defined as the trace of the identity map, that is, the scalar $\mathsf{d}_A: I \to I$ defined as follows: 
\begin{align}
    \mathsf{d}_A = \mathsf{Tr}^A(1_A)
\end{align}
Now in a traced symmetric monoidal closed category, we have the following identity: 
\begin{align}
    \mathsf{d}_A = \mathsf{u}^A;\mathsf{tr}^A
\end{align}
\note{It may be helpful to make the types in the above and such equations explicit --- this was one of the feedback in our earlier paper QPL.}


Now in a dagger traced symmetric monoidal category, we say that $A$ is \textbf{positive-dimensional} if there exists a positive scalar $\mathsf{z}_A: I \to I$ such that: 
\begin{align}
    (\mathsf{z}^2_A \cdot d_A) \cdot 1_A = 1_A 
\end{align}
Now observe that $\mathsf{z}_A \cdot 1_A$ is positive definite with inverse $(\mathsf{z}_A \cdot 1_A)^{-1} = (\mathsf{z}_A \cdot d_A) \cdot 1_A$. Moreover also note that we have the following equality: 
\begin{align}
    \mathsf{z}_A^2 \cdot d_A^2 = d_A
\end{align}
So in the special case that $d_A$ is invertible, we then have that: 
\begin{align}
    \mathsf{z}^2_A = d^{-1}_A
\end{align}


\begin{definition} For a scalar $a$, the \textbf{depolarizing map} at $a$ on $A$ is the map $\mathcal{D}_a:[A,A] \to [A,A]$ defined as follows: 
\begin{align}
    \mathcal{D}_a = a \cdot 1_A + \left( (1_I - a) \cdot \mathsf{z}_A^2 \right) \cdot \left( \mathsf{tr}^A; \mathsf{u}^A\right)
\end{align}
\end{definition}

\begin{lemma} Depolarizing maps are self-adjoint. 
\end{lemma}

\begin{lemma} $\mathcal{D}_a;\mathcal{D}_b = \mathcal{D}_{a \cdot b}$
\end{lemma}

\begin{lemma} If $a$ is MP, then $\mathcal{D}_a$ is MP where $\mathcal{D}^\circ_a = \mathcal{D}_{a^\circ}$. 
\end{lemma}

\begin{lemma} Depolarizing maps are UCPTP.     
\end{lemma}

\begin{lemma} $\mathcal{D}_0$ is a UCPTP which is its own Drazin/$\dagger$-Drazin/MP inverse.     
\end{lemma}

\section{Normalized Dagger Frobenius Algebras}

In this background section, if only to introduce notation, we quickly review normalized dagger Frobenius algebras in a dagger compact closed categories. For an in-depth introduction to dagger compact closed categories, we refer the reader to \cite{ }, and to normalized dagger Frobenius algebras to \cite{ }. 


\subsection{Dagger Compact Closed Categories}



Now recall that a \textbf{compact closed category} is a symmetric monoidal category such that every object has a chosen dual. Explicitly, for every object $A$, there is an object $A^\ast$, called the \textbf{dual} of $A$, equipped with a pair of maps:
\begin{align*}
    \underset{\textbf{cup/counit}}{\cup: A^\ast \otimes A \to I} && \underset{\textbf{cap/unit}}{\cap: I \to A \otimes A^\ast} 
\end{align*}
such that the famous snake equations hold: 
\begin{align}
    (\cap \otimes 1_A);(1_A \otimes \cup) = 1_A && (1_{A^\ast} \otimes \cap);(\cup \otimes 1_{A^\ast}) = 1_{A^\ast}
\end{align}
For convenience, we will work with the convention that $I^\ast = I$ and $(A \otimes B)^\ast = B^\ast \otimes A^\ast$. Now we have an induced contravariant functor $(-)^\ast: \mathbb{X} \to \mathbb{X}$, called the dual functor, which sends objects $A$ to its chosen dual $A^\ast$, and sends a map $f: A \to B$ to the map $f: B^\ast \to A^\ast$ defined as the following composite: 
\begin{align}
\xymatrixcolsep{3pc}\xymatrix{ f^\ast: B^\ast  \ar[r]^-{1 \otimes \cap} & B^\ast \otimes A \otimes A^\ast \ar[r]^-{1 \otimes f \otimes 1} & B^\ast \otimes B \otimes A^\ast \ar[r]^-{\cup \otimes 1} & A^\ast
  } 
\end{align}
Then the cup and cap are extranatural with respect to this functor, that is, the following equalities hold: 
\begin{align*}
    (f^\ast \otimes 1);\cup = (1 \otimes f);\cup && \cap; (f \otimes 1) = \cap;(1 \otimes f^\ast) 
\end{align*}


Then a
By combining the dagger and the dual, we obtain a new (contravariant) functor $(-)_\ast: \mathbb{X} \to \mathbb{X}$ which sends an object $A$ to $A^\ast$ and sends a map $f: A \to B$ to the map $f_\ast: A^\ast \to B^\ast$ defined as follows: 
\begin{align}
    f_\ast = (f^\ast)^\dagger = (f^\dagger)^\ast 
\end{align}
where it is known that the $\dagger$ and $\ast$ commute. It may be useful to record the following identities: 
\begin{equation}\begin{gathered}
    (f_\ast \otimes 1); \cup = (1 \otimes f^\dagger);\cup \qquad \cap;(f^\dagger \otimes 1) = \cap;(1 \otimes f_\ast) \\
    \cup^\dagger;(f_\ast \otimes 1) = \cap^\dagger;(1 \otimes f^\dagger) \qquad (1 \otimes f_\ast);\cap^\dagger = (f^\dagger \otimes 1);\cap^\dagger
\end{gathered}\end{equation}

Lastly, as is done in \cite{}, we will assume that objects in our dagger compact closed categories have \textit{positive dimension}. To do so, we first need to discuss scalars. Recall that in a symmetric monoidal category, a \textbf{scalar} is an endomap on the monoidal unit $a: I \to I$. Now recall that we can always ``scalar multiply'' a map $f: A \to B$ by a scalar $a: I \to I$ to get the map $a \cdot f: A \to B$ defined as $a \cdot f = a \otimes f$. Notice that scalar multiplying a scalar $b$ by another $a$ is simply their composition, $a \cdot b = a;b$, and moreover scalars commute, $a \cdot b = b \cdot a$. Now in a dagger compact closed category, the \textbf{dimension} of an object $A$ is the scalar $d_a$ defined as the following composite: 
\begin{align}
    d_A = \cup^\dagger; \cup = \cap; \cap^\dagger
\end{align}
In other words, the dimension of an object is the \textbf{trace} of its identity map. Now what does it mean for an object to be positive-dimensional? Well first recall that in a dagger category, an endomap $p: A \to A$ is said to be \textbf{positive} if there exists a map $f: A \to B$ such that $p = f;f^\dagger$ (not that $f$ need not be unique). Note that positive maps are always self-adjoint, that is, $p^\dagger= p$. Moreover, also note that by construction, the dimension of an object $A$ is always positive. However we will ask for slightly more. First, we say that an endomap $p: A \to A$ is \textbf{positive definite} if it is both positive and an isomorphism. Then we say that an object $A$ is \textbf{positive dimensional} if there exists a positive definite scalar $\mathsf{z}_A$ such that the following equality holds: 
\begin{align}
    (\mathsf{z}^2_A \cdot d_A) \cdot 1_A = 1_A 
\end{align}
Now observe that $\mathsf{z}_A \cdot 1_A$ is positive definite with inverse $(\mathsf{z}_A \cdot 1_A)^{-1} = (\mathsf{z}_A \cdot d_A) \cdot 1_A$. Moreover also note that we have the following equality: 
\begin{align}
    \mathsf{z}_A^2 \cdot d_A^2 = d_A
\end{align}
So in the special case that $d_A$ is invertible, we then have that: 
\begin{align}
    \mathsf{z}^2_A = d^{-1}_A
\end{align}

For the remainder of the paper (unless otherwise stated), we assume that we are working in a dagger compact closed category such that every object is positive dimensional. Here are now are main two examples of dagger compact closed categories. 

\begin{example} Let $\mathbb{C}$ be the field of complex numbers. Then $\mathsf{FHILB}$ be the category of finite dimensional (complex) Hilbert spaces. 
\end{example}

\begin{example} Let $\mathsf{REL}$ be the category of sets and relations, where recall that in this category, a map $R: X \to Y$ is a subset $R \subseteq X \times Y$. 
\end{example}

\subsection{Normalized Dagger Frobenius Algebras}

Now recall that a \textbf{Frobenius algebra} in a symmetric monoidal category is a tuple $\left( A, \nabla, \Delta, \mathsf{u}, \mathsf{e} \right)$ consisting of an object $A$ and four maps: 
\begin{align*}
    \underset{\textbf{multiplication}}{\nabla: A \otimes A \to A} && \underset{\textbf{comultiplication}}{\Delta: A \to A \otimes A} && \underset{\textbf{unit}}{\mathsf{u}: I \to A} && \underset{\textbf{counit}}{\mathsf{e}: A \to I}
\end{align*}
such that the following equalities hold: 
\begin{gather}
(\nabla \otimes 1);\nabla = (1 \otimes \nabla);\nabla \qquad (1 \otimes \mathsf{u});\nabla = 1 = (\mathsf{u} \otimes 1);\nabla \\
\Delta;(\Delta \otimes 1) = \Delta; (1 \otimes \Delta) \qquad \Delta;(1 \otimes \mathsf{e}) = 1 = \Delta; (\mathsf{e} \otimes 1) \\
(\Delta \otimes 1);(1 \otimes \nabla) = \nabla;\Delta= (1 \otimes \Delta);(\nabla \otimes 1)
\end{gather}
The first line of identities says that $(A, \nabla, \mathsf{u})$ is monoid, the second line says that $(A, \Delta, \mathsf{e})$ is a comonoid, while the third line is the famous Frobenius identities. Now a \textbf{dagger Frobenius algebra} is a Frobenius algebra $\left( A, \nabla, \Delta, \mathsf{u}, \mathsf{e} \right)$ such that the monoid structure and comonoid structure are adjoints of each other, that is, the following equality holds:
\begin{align}
    \nabla^\dagger = \Delta && \mathsf{u}^\dagger = \mathsf{e}
\end{align}
For the definition of being normalized, it will be useful to define the map $\mathsf{e}^\flat: A \to I$ as the partial trace of the multiplication, that is, as the following composite: 
\begin{align}
\xymatrixcolsep{3pc}\xymatrix{ \mathsf{e}^\flat: A \ar[r]^-{1 \otimes \cap} & A \otimes A \otimes A^\ast \ar[r]^-{\nabla \otimes 1} & A \otimes A^\ast  \ar[r]^-{\cap^\dagger} & I
  }  
\end{align}
Then a \textbf{normalizer} for a dagger Frobenius algebra $\underline{A}$ is a positive definite map $\square: A \to A$ which is \textbf{central}, that is, the following equality holds:
\begin{align}
    (\square \otimes 1); \nabla = \nabla;\square = (1 \otimes \square);\nabla && \Delta;(1 \otimes \square) = \square;\Delta = \Delta; (\square \otimes 1)
\end{align}
and also the identity holds: 
\begin{align}
    \square^2; \mathsf{e}^\flat = \mathsf{e} 
\end{align}
Then a \textbf{normalized dagger Frobenius algebra} (or \textbf{N$\dagger$FA} for short) is tuple $\underline{A} = (A, \nabla, \Delta, \mathsf{u}, \mathsf{e}, \square)$ consisting of a dagger Frobenius algebra $(A, \nabla, \Delta, \mathsf{u}, \mathsf{e})$ equipped with a normalizer $\square$. 

Lastly, we will impose an extra condition which did not appear \cite{} but seems natural to assume. We say that a N$\dagger$FA $\underline{A}$ is \textbf{dimension stable} if the following equality holds: 
\begin{align}
    \mathsf{u}; \mathsf{e} = d_A
\end{align}

\begin{example} In \textbf{FHILB}, normalizable dagger Frobenius algebras correspond precisely to finite-dimensional $C^\ast$-algebras (and these are all dimension stable). Of particular interest to the story of the paper is $\mathcal{L}(H)$. 
\begin{align*}
\begin{array}[c]{c}
   \nabla: \mathcal{L}(H) \otimes \mathcal{L}(H) \to \mathcal{L}(H)   \\
     \phi \otimes \varphi \longmapsto \phi \circ \varphi  
\end{array} &&  \begin{array}[c]{c}
\mathsf{u}: \mathbb{C} \longrightarrow \mathcal{L}(H)   \\
     1 \longmapsto \mathsf{id}_H
\end{array}
\end{align*}

If $\lbrace \vert e_1 \rangle, \hdots, \lbrace \vert e_n \rangle \rbrace$ is a basis for $\mathcal{H}$, then $\lbrace \vert e_i \rangle \langle e_j \vert ~\vert~ 1 \leq i,j \leq j \rbrace$ is a basis for $\mathcal{L}(H)$.
\begin{align*}
\begin{array}[c]{c}
   \Delta: \mathcal{L}(H)  \to \mathcal{L}(H) \otimes \mathcal{L}(H)   \\
     \phi = \sum\limits_{i,j} a_{i,j} \vert e_i \rangle \langle e_j \vert \longmapsto \sum\limits_{i,j} \sum\limits_{k} a_{i,j} \vert e_i \rangle \langle e_k \vert \otimes \vert e_k \rangle \langle e_j \vert
\end{array} &&  \begin{array}[c]{c}
\mathsf{e}:  \mathcal{L}(H) \longrightarrow \mathbb{C}    \\
      \phi = \sum\limits_{i,j} a_{i,j} \vert e_i \rangle \langle e_j \vert  \longmapsto \mathsf{Tr}(\phi) = \sum\limits_{i,j} a_{i,i}
\end{array}
\end{align*}

Lastly, the normalizer scalar multiplies a linear map being the inverse 
\begin{align*}
    \square: \mathcal{L}(H) \longrightarrow \mathcal{L}(H) \\
    \phi \longmapsto \frac{1}{\sqrt{\mathsf{dim}(H)}} \cdot \phi
\end{align*}
\end{example}

\begin{example} In $\mathsf{REL}$, normalizable dagger Frobenius algebras correspond precisely to (small) groupoids (and these are dimension stable). As a simple example, let $G$ be a group (in the usual classical sense and written additively) and $\phi: G \to G$ a bijection. Then $G$ is a N$\dagger$FA in $\mathsf{REL}$ where:
\begin{align*}
\begin{array}[c]{c}
   \nabla \subseteq (G \times G) \times G   \\
     \nabla = \lbrace ((g,h), gh) \vert~ \forall g,h \in G \rbrace
\end{array} &&  \begin{array}[c]{c}
\mathsf{u} \subseteq \lbrace \ast \rbrace \times G   \\
    \mathsf{u} = \lbrace (\ast, 1) \rbrace 
\end{array} && \begin{array}[c]{c}
   \Delta \subseteq  G \times (G \times G)   \\
     \Delta = \lbrace (gh, (g,h)) \vert~ \forall g,h \in G \rbrace
\end{array} && \begin{array}[c]{c}
\mathsf{u} \subseteq \lbrace \ast \rbrace \times G   \\
    \mathsf{e} = \lbrace (1,\ast) \rbrace 
\end{array} && \begin{array}[c]{c}
\square_\phi \subseteq G \times G   \\
    \square_\phi = \lbrace (g,\phi(g)) \vert~ \forall g \in G \rbrace 
\end{array}
    \end{align*}
\end{example}

A special class of N$\dagger$FA come from the so called ``pants monoids''. So for every $A$, consider the object $A^\ast \otimes A$. Then $\underline{A^\ast \otimes A} = \left(A^\ast \otimes A, 1_{A^\ast} \otimes \cap^\dagger \otimes 1_A , 1_{A^\ast} \otimes \cap \otimes 1_A, \cap^\dagger, \cup, \mathsf{z}_A \cdot 1_{A^\ast \otimes A} \right)$ is a N$\dagger$FA which moreover dimension stable. 

\section{Quantum Maps}

\subsection{Completely Positive}

It will be useful to introduce some notation. So given a map $g: A \to X \otimes B$, define the map $\langle g \rangle: A^\ast \otimes A \to B^\ast \otimes B$ as follows: 
\begin{align}
\xymatrixcolsep{5pc}\xymatrix{ \langle g \rangle: A^\ast \otimes A \ar[r]^-{g_\ast \otimes g} & B^\ast \otimes X^\ast \otimes X \otimes B \ar[r]^-{1 \otimes \cup \otimes 1} & B^\ast \otimes B
  }  
\end{align}
Notice that for a map $g: A \to B$, we have that $\langle g \rangle = g_\ast \otimes g$. 

Now for a N$\dagger$FA $\underline{A}$, we have canonical maps of $A^\ast$ acting and coacting on $A$: 
\begin{align*}
    \underset{\textbf{action}}{\overline{\nabla}: A^\ast \otimes A \to A} && \underset{\textbf{coaction}}{\overline{\Delta}: A \to A^\ast \otimes A} 
\end{align*}
defined respectively as follows: 
\begin{align}
\xymatrixcolsep{3pc}\xymatrix{ \overline{\nabla}: A^\ast \otimes A \ar[r]^-{1 \otimes \Delta} & A^\ast \otimes A \otimes A \ar[r]^-{\cup \otimes 1} & A
  }  && \xymatrixcolsep{3pc}\xymatrix{ \overline{\Delta}: A \ar[r]^-{\cup^\dagger \otimes 1} & A^\ast \otimes A \otimes A \ar[r]^-{1 \otimes \nabla} & A^\ast \otimes A
  } 
\end{align}
Then for a map $h: A \to X \otimes B$, define the map $[h]: A \to B$ as follows: 
\begin{align}
\xymatrixcolsep{5pc}\xymatrix{ [h]:A \ar[r]^-{\overline{\Delta}} & A^\ast \otimes A \ar[r]^-{\langle g \rangle} & B^\ast \otimes B \ar[r]^-{\overline{\nabla}} & B
  }  
\end{align}

\begin{definition} For N$\dagger$FA $\underline{A}$ and $\underline{B}$, a map $f: A \to B$ is said to be \textbf{completely positive} (or \textbf{CP*} for short) if any of the following equivalent conditions hold: 
\begin{enumerate}[(i)]
\item There exists a map $g: A \to X \otimes B$ such that the following equality holds: 
\begin{align}
    \overline{\nabla}; f; \overline{\Delta} = \langle  g \rangle 
\end{align}
\item There exists a map $h: A \to X \otimes B$ such that the following equality holds: 
\begin{align}
    f = [h]
\end{align}
\end{enumerate}  
\end{definition}

\begin{definition} A map $f: A^\ast \otimes A \to B^\ast \otimes B$ is said to be \textbf{completely positive} (or \textbf{CP} for short) if there exists a map $g: A \to X \otimes B$ such that $f = \langle g \rangle$. 
\end{definition}

\begin{theorem} A map $f: A^\ast \otimes A \to B^\ast \otimes B$ is CP if and only if it is CP*
\end{theorem}


\subsection{Trace Preserving}

\begin{definition} For N$\dagger$FA $\underline{A}$ and $\underline{B}$, a map $f: A \to B$ is said to be \textbf{trace preserving} (or \textbf{TP*} for short) if the following equality holds:
\begin{align}
f; \mathsf{e} = \mathsf{e}
\end{align}
\end{definition}

\begin{definition} A map $f: A^\ast \otimes A \to B^\ast \otimes B$ is said to be \textbf{trace preserving} (\textbf{TP} for short) if the following equality holds: 
\begin{align}
f; \cup = \cup
\end{align}
\end{definition}

\begin{example} Here is an example of a TP map whose adjoint is not TP.     
\end{example}

\begin{definition}
    
\end{definition}



\subsection{Quantum Maps}

\begin{definition} For N$\dagger$FA $\underline{A}$ and $\underline{B}$, a \textbf{quantum map} (or \textbf{CPTP*} for short) is a map $f: A \to B$ which is CP* and TP*. 
\end{definition}

These were called normalized CP* maps in \cite{}. 

\begin{lemma} For N$\dagger$FA $\underline{A}$ and $\underline{B}$, a CP* map $f: A \to B$ is TP* if and only if there exists a map $h: A \to X \otimes B$ such that $f = [h]$ and the following equality holds: 
\begin{align}
    \overline{\Delta}; \langle h \rangle; \cup = \mathsf{e}
\end{align}
\end{lemma}

\begin{corollary} For N$\dagger$FA $\underline{A}$ and $\underline{B}$, if for a CP* map $f: A \to B$, there exists a map $h: A \to X \otimes B$ such that $f = [h]$ and also satisfies:
\begin{align}
    \langle h \rangle; \cup = (1 \otimes \square^2); \cup
\end{align}
Then $f$ is TP*. 
\end{corollary}

\begin{definition} For N$\dagger$FA $\underline{A}$ and $\underline{B}$, a \textbf{dagger quantum map} (or \textbf{CP$\dagger$TP*} for short) is a map $f: A \to B$ which is CP* and $\dagger$TP*. 
\end{definition}

\subsection{Unital}



\section{Some Examples of CPTP Maps}




\section{Trace Preserving Maps}

\begin{definition} Let $\underline{A}$ and $\underline{B}$ be normalized dagger Frobenius algebras. Then a map $f: A \to B$ is said to be: 
\begin{enumerate}[(i)]
\item \textbf{trace preserving} (\textbf{TP*} for short) if it preserves the counit, that is, if the following equality holds: 
\begin{align}
f; \mathsf{e} = \mathsf{e}
\end{align}
    \item \textbf{cotrace preserving} (\textbf{coTP*} for short) if it preserves the unit, that is, if the following equality holds: 
    \begin{align}
\mathsf{u}; f = \mathsf{u}
\end{align}
\end{enumerate}     
\end{definition}

\begin{lemma} Let $\underline{A}$ and $\underline{B}$ be normalized dagger Frobenius algebras. Then for a map $f: A \to B$, $f$ is TP* (resp. coTP*) if and only if $f^\dagger$ is coTP* (resp. TP*).  
\end{lemma}

\begin{definition} A map $f: A^\ast \otimes A \to B^\ast \otimes B$ is said to be: 
\begin{enumerate}[(i)]
\item \textbf{trace preserving} (\textbf{TP} for short) if the following equality holds: 
\begin{align}
f; \cup = \cup
\end{align}
    \item \textbf{cotrace preserving} (\textbf{coTP} for short) if the following equality holds: 
    \begin{align}
\cup^\dagger; f = \cap; \sigma; f = \cap; \sigma = \cup^\dagger
\end{align}
\end{enumerate}  
\end{definition}

Here is an example of TP map whose adjoints is not TP. 

\begin{example}
    
\end{example}

\begin{definition} Let $\underline{A}$ and $\underline{B}$ be normalized dagger Frobenius algebras. Then a map $f: A \to B$ is said to be \textbf{$\dagger$-trace preserving} (\textbf{$\dagger$-TP*} for short) if both $f$ and $f^\dagger$ are TP*. 
\end{definition}

\begin{lemma} Let $\underline{A}$ and $\underline{B}$ be normalized dagger Frobenius algebras. Then a map $f: A \to B$, the following are equivalent: 
\begin{enumerate}[(i)]
\item $f$ is $\dagger$-TP*;
\item $f$ is TP* and coTP*;
\item $f$ and $f^\dagger$ are coTP*. 
\end{enumerate}  
\end{lemma}
\begin{proof}
\end{proof}

\section{Drazin Inverses}

\begin{proposition} The Drazin inverse of a (co)TP* endomorphism is (co)TP*. 
\end{proposition}


\section{Dagger-Drazin Inverses}

\begin{proposition} Let $\underline{A}$ and $\underline{B}$ be normalized dagger Frobenius algebras. If $f :A \to B$ is $\dagger$-Drazin and $\dagger$-TP*, then $f^\partial: B \to A$ is $\dagger$-TP*. 
\end{proposition}

\section{Moore-Penrose Inverses}

\begin{proposition} Let $\underline{A}$ and $\underline{B}$ be normalized dagger Frobenius algebras. If $f :A \to B$ is MP then $f$ is $\dagger$-TP* if and only if $f^\circ$ is $\dagger$-TP*. 
\end{proposition}

\section{Quantum Maps}

\begin{definition} In a $\dagger$-compact closed category, a map $\Phi: {[A,A]} \to {[B,B]}$ is \textbf{completely positive} (or \textbf{CP} for short) if its associated transpose $\Phi^\sharp: A \otimes B \to A \otimes B$ is positive. 
\end{definition}


\begin{definition} In a $\dagger$-compact closed category, a \textbf{quantum map} (or CPTP short) is a map $\Phi: {[A,A]} \to {[B,B]}$ which is both CP and TP. 
\end{definition}

Important, quantum maps need not be $\dagger$-TP. As such, the MP inverse (and hence a $\dagger$-Drazin inverse) of a quantum map need not be a quantum map either Here is an example. 

\section{Pure Unitary Maps}

\begin{lemma} In a compact closed category, let $u: A \to B$ and $v: B \to A$ be a section-retraction pair (i.e. $u;v = 1_A$). Then $v^\ast \otimes u: A^\ast \otimes A \to B^\ast \otimes B$ is TP. 
\end{lemma}
\begin{proof} 
\begin{align*}
   (v^\ast \otimes u); \epsilon_B = \left( 1_{A^\ast} \otimes (u;v) \right); \epsilon_B = (1_{A^\ast} \otimes 1_A); \epsilon_B = \epsilon_B
\end{align*}
\end{proof}

\begin{corollary} In a dagger compact closed category, let $u: A \to B$ be a unitary map. Then $u_\ast \otimes u: A^\ast \otimes A \to B^\ast \otimes B$ is $\dagger$-TP.
\end{corollary}

\section{Mixed Unitary Maps}

\begin{definition} In a $\dagger$-additive symmetric monoidal category, a \textbf{partition of unity} is a finite family of scalars $\lbrace a_i: I \to I \rbrace^n_{i=1}$ such each $a_i$ is positive and $\sum \limits_i a_i = 1_I$. 
\end{definition}

\begin{definition} In a $\dagger$-additive traced symmetric monoidal closed category, a \textbf{mixed unitary map} is a map $\Phi: A^\ast \otimes A \to B^\ast \otimes B$ such that there exists a finite family of unitary maps $\lbrace u_i: X \to Y \vert~ \text{$u_i$ unitary} \rbrace^n_{i=1}$ and a partition of unity $\lbrace a_i: I \to I \rbrace^n_{i=1}$ such that the following equality holds:
\begin{align}
    \Phi = \sum\limits^n_{i=1} a_i \cdot \left( u_\ast \otimes u \right)
\end{align}
\end{definition}

\begin{lemma} Every mixed unitary map is $\dagger$-TP. 
\end{lemma}

\begin{corollary} If a mixed unitary map $\Phi$ is MP, then $\Phi^\circ$ is $\dagger$-TP. 
\end{corollary}

\begin{proposition} In an $\dagger$-additive compact closed category, every mixed unitary map is CPTP. 
\end{proposition}
\begin{proof}
    
\end{proof}

\section{Depolarizing Maps}

\begin{definition} In a dagger compact closed category, an object $A$ is said to be \textbf{positive dimensional} if there exists a positive scalar $z: A \to A$ such that the following equality holds: 
\begin{align}
    \left( z \cdot \mathsf{dim}(A) \right) \cdot 1_A = 1_A
\end{align}
\end{definition}

\begin{definition} Let $\underline{A}$ be a normalized dagger Frobenius algebra. For every scalar $a: I \to I$, define its associated \textbf{depolarizing map} on $A$ as the map $\mathcal{D}_a: A \to A$ defined as follows: 
\begin{align}
    \mathcal{D}_a = a \cdot 1_A + \left( (1_I - a) \cdot \mathsf{z}_A^2 \right) \cdot \left( \mathsf{u}; \mathsf{e}; \square^{-1}; \square^{-1} \right)
\end{align}
    
\end{definition}

\begin{definition} In a dagger compact closed category, let $a: I \to I$ be a scalar. The \textbf{depolarizing map} at $a$ on $A$ is the map $\mathcal{D}_a: A \otimes A^\ast \to A^\ast \otimes A$ defined as follows: 
\begin{align*}
   \mathcal{D}_a = a \cdot 1_{A \otimes A^\ast} + (1_I - a) \cdot \mathsf{dim}(A)^{-1} \cdot \left( \epsilon_A; \epsilon^\dagger_A \right)
\end{align*}
\end{definition}

\begin{lemma} Every depolarizing map is $\dagger$-TP. 
\end{lemma}



\begin{proposition} Let $A$ be of positive dimensional, if $a$ and $1-a$ are positive, then $\mathcal{D}_a$ is a quantum map. 
\end{proposition} 
